\documentclass[11pt,a4paper]{article}

% ==================== TEMEL PAKETLER (pdflatex) ====================
\usepackage[T1]{fontenc}
\usepackage[utf8]{inputenc}
\usepackage{lmodern}
\usepackage[a4paper,margin=2cm,top=2.4cm,bottom=2.4cm]{geometry}
\usepackage{xcolor}
\usepackage{colortbl}
\usepackage{array}
\usepackage{booktabs}
\usepackage{longtable}
\usepackage{tabularx}
\usepackage{listings}
\usepackage{fancyhdr}
\usepackage{amsmath}
\usepackage{amssymb}
\usepackage{bm}
\usepackage{hyperref}

% ==================== RENK PALETİ ====================
\definecolor{anaMavi}{HTML}{1B3A5B}
\definecolor{vurguMavi}{HTML}{2E6DA4}
\definecolor{acikMavi}{HTML}{EAF2F9}
\definecolor{turuncu}{HTML}{D9701E}
\definecolor{yesil}{HTML}{2E7D32}
\definecolor{acikYesil}{HTML}{E8F5E9}
\definecolor{kirmizi}{HTML}{C0392B}
\definecolor{mor}{HTML}{6C3483}
\definecolor{acikMor}{HTML}{F3E9F7}
\definecolor{acikGri}{HTML}{F4F5F7}
\definecolor{ortaGri}{HTML}{6B7280}
\definecolor{koyuGri}{HTML}{2D2D2D}
\definecolor{kodArka}{HTML}{FBFBFA}
\definecolor{formulArka}{HTML}{F7F9FC}

% ==================== HYPERREF ====================
\hypersetup{
  colorlinks=true,
  linkcolor={[HTML]{2E6DA4}},
  urlcolor={[HTML]{2E6DA4}},
  pdftitle={Makine Öğrenmesi: Matematikten Koda Kapsamlı Rehber},
  pdfauthor={ML Rehberi},
  bookmarksnumbered=true
}

% ==================== BAŞLIK STİLLERİ (titlesec olmadan) ====================
\setcounter{secnumdepth}{0}
\newcounter{secc}
\newcounter{subsecc}[secc]
\newcounter{subsubsecc}[subsecc]
\newcommand{\gsec}[1]{\par\phantomsection\stepcounter{secc}%
  \setcounter{subsecc}{0}%
  \vspace{18pt}{\color{anaMavi}\Large\bfseries\thesecc.\hspace{0.5em}#1}\par
  \vspace{2pt}{\color{turuncu}\rule{\linewidth}{1.6pt}}\par\vspace{8pt}%
  \addcontentsline{toc}{section}{\protect\numberline{\thesecc}#1}\markboth{#1}{#1}}
\newcommand{\gsub}[1]{\par\phantomsection\stepcounter{subsecc}%
  \setcounter{subsubsecc}{0}%
  \vspace{11pt}{\color{vurguMavi}\large\bfseries\thesecc.\thesubsecc\hspace{0.5em}#1}\par\vspace{4pt}%
  \addcontentsline{toc}{subsection}{\protect\numberline{\thesecc.\thesubsecc}#1}}
\newcommand{\gsubsub}[1]{\par\phantomsection\stepcounter{subsubsecc}%
  \vspace{7pt}{\color{koyuGri}\normalsize\bfseries #1}\par\vspace{2pt}}

% Kısım (Part) başlığı
\newcommand{\gpart}[2]{\clearpage
  \begin{center}\vspace*{1.5cm}
  {\color{turuncu}\rule{0.7\linewidth}{1.5pt}}\\[10pt]
  {\color{ortaGri}\Large KISIM #1}\\[8pt]
  {\color{anaMavi}\Huge\bfseries #2}\\[10pt]
  {\color{turuncu}\rule{0.7\linewidth}{1.5pt}}
  \end{center}\vspace{1cm}}

% ==================== BAŞLIK / ALTBİLGİ ====================
\setlength{\headheight}{14pt}
\pagestyle{fancy}
\fancyhf{}
\renewcommand{\headrulewidth}{0.4pt}
\renewcommand{\footrulewidth}{0.4pt}
\fancyhead[L]{\small\color{ortaGri}Makine Öğrenmesi Rehberi}
\fancyhead[R]{\small\color{ortaGri}\leftmark}
\fancyfoot[C]{\small\color{ortaGri}\thepage}

% ==================== KOD LİSTELEME (Python) ====================
\lstdefinestyle{pystyle}{
  language=Python,
  basicstyle=\ttfamily\small\color{koyuGri},
  keywordstyle=\color{vurguMavi}\bfseries,
  commentstyle=\color{yesil}\itshape,
  stringstyle=\color{turuncu},
  numberstyle=\tiny\color{ortaGri},
  numbers=left,
  numbersep=8pt,
  backgroundcolor=\color{kodArka},
  frame=single,
  framesep=6pt,
  rulecolor=\color{ortaGri!40},
  breaklines=true,
  showstringspaces=false,
  tabsize=4,
  columns=fullflexible,
  keepspaces=true,
  extendedchars=true,
  literate={ı}{{\i}}1{İ}{{\.I}}1{ş}{{\c s}}1{Ş}{{\c S}}1{ğ}{{\u g}}1{Ğ}{{\u G}}1{ç}{{\c c}}1{Ç}{{\c C}}1{ö}{{\"o}}1{Ö}{{\"O}}1{ü}{{\"u}}1{Ü}{{\"U}}1,
  morekeywords={np,pd,plt,sns,as,with,yield,lambda,None,True,False,
    self,assert,async,await,nonlocal,fit,transform,predict,fit_transform}
}
\lstset{style=pystyle}

% Kabuk/terminal stili
\lstdefinestyle{shstyle}{
  basicstyle=\ttfamily\small\color{koyuGri},
  backgroundcolor=\color{koyuGri!6},
  frame=single,framesep=6pt,rulecolor=\color{ortaGri!40},
  breaklines=true,showstringspaces=false,tabsize=4,columns=fullflexible,
  keepspaces=true,
  literate={ı}{{\i}}1{İ}{{\.I}}1{ş}{{\c s}}1{Ş}{{\c S}}1{ğ}{{\u g}}1{Ğ}{{\u G}}1{ç}{{\c c}}1{Ç}{{\c C}}1{ö}{{\"o}}1{Ö}{{\"O}}1{ü}{{\"u}}1{Ü}{{\"U}}1,
}

% ==================== ÖZEL KUTULAR (tcolorbox olmadan) ====================
\newcommand{\callout}[4]{\par\smallskip\noindent\begingroup
  \setlength{\fboxsep}{8pt}\setlength{\fboxrule}{0.8pt}%
  \fcolorbox{#1}{#2}{\parbox{\dimexpr\linewidth-2\fboxsep-2\fboxrule\relax}{%
    {\bfseries\color{#1}#3}\par\smallskip #4}}\endgroup\par\smallskip}
\newcommand{\bilgi}[1]{\callout{vurguMavi}{acikMavi}{Bilgi}{#1}}
\newcommand{\ipucu}[1]{\callout{yesil}{acikYesil}{İpucu}{#1}}
\newcommand{\uyari}[1]{\callout{kirmizi}{kirmizi!7}{Dikkat}{#1}}
\newcommand{\ornek}[2][Örnek]{\callout{ortaGri!90!black}{acikGri}{#1}{#2}}
\newcommand{\notk}[2][Not]{\callout{mor}{acikMor}{#1}{#2}}
% Matematik/turetme kutusu
\newcommand{\turetme}[1]{\callout{anaMavi}{formulArka}{Türetme}{#1}}

% Yardımcı komutlar
\newcommand{\code}[1]{\texttt{\color{koyuGri}#1}}
\newcommand{\vek}[1]{\bm{#1}}
\newcommand{\R}{\mathbb{R}}
\renewcommand{\arraystretch}{1.35}

% ==================== BAŞLANGIÇ ====================
\begin{document}

% ---------- KAPAK ----------
\begin{titlepage}
\centering
\vspace*{1.4cm}
{\color{turuncu}\rule{\textwidth}{2pt}}\\[8pt]
{\color{anaMavi}\Huge\bfseries Makine Öğrenmesi}\\[10pt]
{\color{vurguMavi}\LARGE Matematikten Koda Kapsamlı Rehber}\\[8pt]
{\color{ortaGri}\large Teori \textbullet\ Türetmeler \textbullet\ Sıfırdan ve scikit-learn ile Uygulama}\\[6pt]
{\color{turuncu}\rule{\textwidth}{2pt}}\\[1.4cm]

\begin{center}
\fcolorbox{vurguMavi}{acikMavi}{\parbox{0.86\textwidth}{\centering\large
Bu belge, makine öğrenmesini \textbf{matematiksel temellerinden başlayarak}
her algoritmanın türetmesiyle, hem \textbf{sıfırdan NumPy} hem de
\textbf{scikit-learn} uygulamalarıyla adım adım anlatır. Lineer cebir ve
optimizasyondan derin öğrenmeye ve geri yayılıma kadar uzanır.\vspace{4pt}}}
\end{center}

\vspace{0.7cm}
\begin{center}
\fcolorbox{ortaGri!50}{acikGri}{\parbox{0.86\textwidth}{\small
\textbf{Kısım 1 --- Matematiksel Temeller:} lineer cebir \textbullet\ kalkülüs
\& gradyanlar \textbullet\ optimizasyon (gradient descent) \textbullet\ olasılık
\& MLE \textbullet\ bilgi teorisi.\vspace{3pt}

\textbf{Kısım 2 --- Öğrenme Teorisi:} kayıp fonksiyonları \textbullet\ risk
\textbullet\ bias-variance \textbullet\ düzenlileştirme \textbullet\ değerlendirme.\vspace{3pt}

\textbf{Kısım 3 --- Regresyon:} lineer regresyon (normal denklem + GD)
\textbullet\ Ridge/Lasso \textbullet\ lojistik regresyon.\vspace{3pt}

\textbf{Kısım 4 --- Sınıflandırma:} Naive Bayes \textbullet\ KNN \textbullet\
SVM (marj + kernel) \textbullet\ karar ağaçları \textbullet\ ensemble \& boosting.\vspace{3pt}

\textbf{Kısım 5 --- Denetimsiz:} K-Means \textbullet\ PCA (özdeğer türetmesi).\vspace{3pt}

\textbf{Kısım 6 --- Derin Öğrenme:} perceptron \textbullet\ çok katmanlı ağlar
\textbullet\ geri yayılım (tam türetme) \textbullet\ optimizasyon \textbullet\ CNN \& RNN.\vspace{3pt}}}
\end{center}

\vfill
{\color{ortaGri}\large Hazırlanma Tarihi: 12 Temmuz 2026}\\[4pt]
{\color{ortaGri} Python 3 \textbullet\ NumPy \textbullet\ scikit-learn \textbullet\ Kodlar İngilizce, açıklamalar Türkçe}
\vspace*{0.6cm}
\end{titlepage}

% ---------- İÇİNDEKİLER ----------
\tableofcontents
\thispagestyle{empty}
\clearpage

% #####################################################################
\gpart{1}{Matematiksel Temeller}
% #####################################################################

Makine öğrenmesi, özünde \emph{matematiktir}: veriyi vektör ve matrislerle
temsil eder (lineer cebir), bir hata fonksiyonunu türev alarak küçültür
(kalkülüs \& optimizasyon) ve belirsizliği olasılıkla modeller. Bu kısım,
sonraki tüm algoritmaların üzerine kurulacağı matematiksel zemini atar.
Aceleyle geçme --- buradaki her kavram ileride tekrar tekrar karşına çıkacak.

% =====================================================================
\gsec{Lineer Cebir}
% =====================================================================

Lineer cebir, makine öğrenmesinin dilidir. Bir veri kümesindeki her örnek bir
\emph{vektör}, tüm veri kümesi bir \emph{matris}, bir modelin parametreleri de
vektör/matrislerdir. Modern donanım (CPU/GPU) matris işlemlerinde son derece
hızlı olduğundan, algoritmaları matris biçiminde ifade etmek hem zarif hem
verimlidir.

\gsub{Skaler, Vektör, Matris, Tensör}

\begin{center}
\begin{tabularx}{\textwidth}{@{}>{\bfseries\color{anaMavi}}l l X@{}}
\toprule
\rowcolor{acikMavi}\color{anaMavi}\textbf{Nesne} & \color{anaMavi}\textbf{Gösterim} & \color{anaMavi}\textbf{Anlamı}\\
\midrule
Skaler & $x \in \R$ & Tek sayı.\\
\rowcolor{acikGri} Vektör & $\vek{x} \in \R^n$ & $n$ sayıdan oluşan dizi (bir örnek/nokta).\\
Matris & $\vek{X} \in \R^{m\times n}$ & $m$ satır $n$ sütun ($m$ örnek, $n$ özellik).\\
\rowcolor{acikGri} Tensör & $\mathcal{X}$ & 3+ boyutlu dizi (görüntü, video).\\
\bottomrule
\end{tabularx}
\end{center}

Bir makine öğrenmesi problemi genelde \emph{tasarım matrisi} $\vek{X} \in
\R^{m \times n}$ ile ifade edilir: her satır bir örnek ($m$ tane), her sütun
bir özellik ($n$ tane). Hedef değerler ise $\vek{y} \in \R^m$ vektöründedir.

\gsub{İç Çarpım (Dot Product)}

İki vektörün iç çarpımı, makine öğrenmesindeki en temel işlemdir: bir modelin
``ağırlıklı toplam'' tahmini tam olarak budur. Tanım:
\[
\vek{a} \cdot \vek{b} = \vek{a}^\top \vek{b} = \sum_{i=1}^{n} a_i b_i
= a_1 b_1 + a_2 b_2 + \dots + a_n b_n .
\]
İç çarpım aynı zamanda vektörler arası açıyla ilişkilidir:
$\vek{a}^\top \vek{b} = \|\vek{a}\|\,\|\vek{b}\|\cos\theta$. Bu yüzden iç
çarpım bir \emph{benzerlik} ölçüsüdür: aynı yöne bakan vektörlerde büyük,
dik vektörlerde sıfırdır.

\begin{lstlisting}
import numpy as np

a = np.array([1.0, 2.0, 3.0])
b = np.array([4.0, 5.0, 6.0])

dot = a @ b                 # 1*4 + 2*5 + 3*6 = 32
dot = np.dot(a, b)          # aynı sonuç
dot = np.sum(a * b)         # elle: eleman bazlı çarpım sonra toplam
\end{lstlisting}

\gsub{Norm: Vektör Uzunluğu}

Norm, bir vektörün ``büyüklüğünü'' ölçer. Kayıp fonksiyonları ve
düzenlileştirme normlarla tanımlanır. En yaygını $L_2$ (Öklid) normudur:
\[
\|\vek{x}\|_2 = \sqrt{\sum_{i=1}^{n} x_i^2} = \sqrt{\vek{x}^\top \vek{x}},
\qquad
\|\vek{x}\|_1 = \sum_{i=1}^{n} |x_i| .
\]
$L_1$ normu (mutlak değerler toplamı) Lasso düzenlileştirmesinde,
$L_2$ normu Ridge'de karşımıza çıkacak. İkisi arasındaki fark, ileride göreceğimiz
gibi, modelin davranışını kökten değiştirir.

\begin{lstlisting}
x = np.array([3.0, 4.0])
l2 = np.linalg.norm(x)          # sqrt(9 + 16) = 5.0
l1 = np.linalg.norm(x, ord=1)   # |3| + |4| = 7.0
\end{lstlisting}

\gsub{Kosinüs Benzerliği ve Uzaklık Ölçütleri}

İki veri noktasının ``ne kadar benzer'' veya ``ne kadar uzak'' olduğunu ölçmek;
KNN, k-means, öneri sistemleri ve metin gömme (embedding) gibi birçok
algoritmanın temelidir. İki temel yaklaşım vardır: \emph{açı} tabanlı benzerlik
ve \emph{uzaklık}.

\gsubsub{Kosinüs Benzerliği (Cosine Similarity)}

İç çarpımdan türeyen kosinüs benzerliği, iki vektörün \emph{yönlerinin} ne kadar
hizalı olduğunu ölçer --- büyüklüklerinden bağımsız olarak:
\[
\cos\theta = \frac{\vek{a}^\top\vek{b}}{\|\vek{a}\|\,\|\vek{b}\|} \in [-1, 1].
\]
$+1$ aynı yön (çok benzer), $0$ dik (ilgisiz), $-1$ zıt yön demektir.

\turetme{\textbf{Örnek:} $\vek{a}=(1,2,3)$, $\vek{b}=(2,4,6)$ olsun. Dikkat:
$\vek{b}=2\vek{a}$, yani aynı yön. İç çarpım $\vek{a}^\top\vek{b}=2+8+18=28$;
normlar $\|\vek{a}\|=\sqrt{14}$, $\|\vek{b}\|=\sqrt{56}=2\sqrt{14}$. O halde:
\[
\cos\theta = \frac{28}{\sqrt{14}\cdot 2\sqrt{14}} = \frac{28}{28} = 1.
\]
Büyüklükleri farklı olsa da (biri diğerinin iki katı), yönleri aynı olduğu için
benzerlik tamdır. Buna karşılık $\vek{a}=(1,0)$ ve $\vek{b}=(0,1)$ için iç
çarpım $0$ olduğundan $\cos\theta=0$: dik, yani ``ilgisiz''.}

\gsubsub{Uzaklık Ölçütleri (Distance Metrics)}

Uzaklık ise iki noktanın \emph{konum} farkını ölçer. İki yaygın ölçüt:
\[
\text{Öklid (}L_2\text{): } d_2(\vek{a},\vek{b}) = \|\vek{a}-\vek{b}\|_2
= \sqrt{\textstyle\sum_i (a_i-b_i)^2}, \qquad
\text{Manhattan (}L_1\text{): } d_1 = \sum_i |a_i-b_i|.
\]
Örneğin $\vek{a}=(1,2)$, $\vek{b}=(4,6)$ için: $d_2=\sqrt{9+16}=5$,
$d_1=3+4=7$.

\begin{lstlisting}
import numpy as np

a = np.array([1.0, 2.0, 3.0])
b = np.array([2.0, 4.0, 6.0])

cos_sim = a @ b / (np.linalg.norm(a) * np.linalg.norm(b))   # 1.0 (aynı yön)
euclidean = np.linalg.norm(a - b)                            # Öklid uzaklığı
manhattan = np.sum(np.abs(a - b))                            # Manhattan uzaklığı
\end{lstlisting}

\bilgi{\textbf{Benzerlik mi, uzaklık mı?} Kosinüs benzerliği büyüklüğü
\emph{yok sayar}, yalnızca yöne bakar --- metin/embedding karşılaştırmasında
ideal (uzun bir belge ile kısa bir belge aynı konudaysa benzer sayılmalı).
Öklid uzaklığı ise büyüklüğe \emph{duyarlıdır}; bu yüzden KNN ve k-means'te
özellikleri önce standartlaştırmak (aynı ölçeğe getirmek) şarttır.}

\gsub{Matris Çarpımı}

Matris çarpımı, bir veri kümesindeki tüm örnekler için tahminleri \emph{tek
seferde} hesaplamamızı sağlar. $\vek{A} \in \R^{m\times k}$ ve
$\vek{B} \in \R^{k\times n}$ için çarpım $\vek{C} = \vek{AB} \in \R^{m\times n}$:
\[
C_{ij} = \sum_{l=1}^{k} A_{il} B_{lj} .
\]
Çarpımın tanımlı olması için ilk matrisin sütun sayısı, ikincinin satır
sayısına eşit olmalıdır. Bu ``iç boyutların eşleşmesi'' kuralı,
uygulamada boyut hatalarının ana kaynağıdır.

\begin{lstlisting}
A = np.array([[1, 2],
              [3, 4]])          # 2x2
B = np.array([[5, 6],
              [7, 8]])          # 2x2

C = A @ B                       # matris çarpımı
# C = [[1*5+2*7, 1*6+2*8],
#      [3*5+4*7, 3*6+4*8]] = [[19, 22], [43, 50]]

# TÜM örnekler için tek seferde bir model tahmini:
# X: (m, n) tasarım matrisi, w: (n,) ağırlıklar  ->  tahminler: (m,)
X = np.random.randn(100, 3)     # 100 örnek, 3 özellik
w = np.array([0.5, -1.2, 2.0])
predictions = X @ w             # (100,) -- her satır için bir iç çarpım
\end{lstlisting}

\gsub{Şekil Uyumu ve Yayınlama (Broadcasting)}

Matris/vektör işlemlerinde en sık karşılaşılan hata, \emph{boyut uyumsuzluğudur}
(shape mismatch). İki kuralı içselleştirmek, ML kodundaki hataların büyük
kısmını önler:

\begin{itemize}\itemsep2pt
\item \textbf{Matris çarpımı} $\vek{A}_{(m\times k)}\vek{B}_{(k\times n)}$:
iç boyutlar ($k$) eşleşmeli; sonuç $(m\times n)$.
\item \textbf{Yayınlama (broadcasting):} eleman bazlı işlemlerde NumPy şekilleri
\emph{sondan başa} karşılaştırır; her boyut ya eşit ya da biri $1$ olmalıdır.
Boyutu $1$ olan taraf diğerine ``yayılır'' --- bellekte kopya oluşturmadan.
\end{itemize}

\begin{center}
\begin{tabularx}{\textwidth}{@{}>{\ttfamily}l >{\ttfamily}l >{\ttfamily}l X@{}}
\toprule
\rowcolor{acikMavi}\color{anaMavi}\textbf{A şekli} & \color{anaMavi}\textbf{B şekli} & \color{anaMavi}\textbf{Sonuç} & \color{anaMavi}\textbf{Durum}\\
\midrule
(100, 3) & (3,) & (100,) & A @ B: geçerli (iç boyut 3)\\
\rowcolor{acikGri} (100, 3) & (3,) & (100, 3) & A + B: her satıra yayılır\\
(100, 1) & (1, 3) & (100, 3) & Dış çarpım gibi yayılır\\
\rowcolor{acikGri} (100, 3) & (100,) & --- & Hata: 3 ile 100 uyumsuz\\
\bottomrule
\end{tabularx}
\end{center}

\begin{lstlisting}
import numpy as np

X = np.random.randn(100, 3)     # 100 örnek, 3 özellik

# Her sütunu standartlaştır (z-skoru) -- yayınlamanın klasik kullanımı
mean = X.mean(axis=0)           # (3,)  her özelliğin ortalaması
std = X.std(axis=0)             # (3,)
X_scaled = (X - mean) / std     # (100,3) - (3,) -> (3,) her satıra yayılır

# Sık hata: (100,) ile (100,3) toplanamaz -> yeniden şekillendir
col = np.arange(100)            # (100,)
result = X + col[:, None]       # (100,1) yapılıp her sütuna yayılır
\end{lstlisting}

\uyari{Boyut hatası aldığında (\code{ValueError: operands could not be
broadcast}), önce ilgili dizilerin \code{.shape}'ini yazdır. Sorun neredeyse
her zaman bir yerde beklenmedik bir eksik/fazla boyuttur. \code{None}
(veya \code{np.newaxis}) ile yeni bir eksen ekleyerek şekilleri uyumlu hale
getirebilirsin.}

\gsub{Özel Matrisler ve İşlemler}

\begin{center}
\begin{tabularx}{\textwidth}{@{}>{\bfseries\color{anaMavi}}l X@{}}
\toprule
\rowcolor{acikMavi}\color{anaMavi}\textbf{Kavram} & \color{anaMavi}\textbf{Tanım / önem}\\
\midrule
Transpoz $\vek{A}^\top$ & Satır ve sütunları takas eder; $(A^\top)_{ij}=A_{ji}$.\\
\rowcolor{acikGri} Birim matris $\vek{I}$ & Köşegeni 1, gerisi 0; $\vek{AI}=\vek{A}$.\\
Ters matris $\vek{A}^{-1}$ & $\vek{A}\vek{A}^{-1}=\vek{I}$; sadece kare ve tekil olmayan matrisler için.\\
\rowcolor{acikGri} Determinant $\det(\vek{A})$ & 0 ise matris tersinir değildir (tekil).\\
Simetrik & $\vek{A}=\vek{A}^\top$; kovaryans matrisleri simetriktir.\\
\bottomrule
\end{tabularx}
\end{center}

\gsub{Özdeğerler ve Özvektörler}

Bir kare matris $\vek{A}$ için, yönü değişmeyip yalnızca ölçeklenen özel
vektörlere \emph{özvektör}, ölçekleme çarpanına \emph{özdeğer} denir:
\[
\vek{A}\vek{v} = \lambda \vek{v}, \qquad \vek{v} \ne \vek{0}.
\]
Burada $\vek{v}$ özvektör, $\lambda$ özdeğerdir. Özdeğerler, PCA'nın (Kısım 5)
kalbindedir: bir veri kümesinin kovaryans matrisinin özvektörleri, verinin en
çok yayıldığı yönleri (temel bileşenleri) verir; özdeğerler o yönlerdeki
varyansı ölçer.

\begin{lstlisting}
A = np.array([[2.0, 0.0],
              [0.0, 3.0]])
eigenvalues, eigenvectors = np.linalg.eig(A)
# özdeğerler:   [2., 3.]
# özvektörler: sütunlar özvektörlerdir
\end{lstlisting}

\bilgi{\textbf{Neden matris biçimi?} Bir döngüyle 100 örnek için tahmin
hesaplamak yerine tek bir \code{X @ w} çarpımı kullanmak hem çok daha hızlıdır
(BLAS kütüphaneleri, GPU) hem de matematiği okunaklı kılar. Makine öğrenmesi
kodunu ``vektörleştirmek'' (döngüleri matris işlemleriyle değiştirmek) temel
bir beceridir.}

% =====================================================================
\gsec{Kalkülüs ve Optimizasyon}
% =====================================================================

Makine öğrenmesinde ``öğrenmek'', bir \emph{kayıp fonksiyonunu} (modelin ne
kadar yanıldığını ölçen fonksiyon) minimize etmektir. Bunu yapmak için kalkülüs
kullanırız: türev, fonksiyonun hangi yönde arttığını söyler; ters yöne giderek
kaybı azaltırız. Bu, \emph{gradient descent} (gradyan inişi) algoritmasının özüdür.

\gsub{Türev: Değişim Oranı}

Tek değişkenli bir $f(x)$ fonksiyonunun türevi $f'(x)$, fonksiyonun $x$
noktasındaki anlık değişim oranıdır (eğim):
\[
f'(x) = \frac{df}{dx} = \lim_{h \to 0} \frac{f(x+h) - f(x)}{h}.
\]
Bir minimum (veya maksimum) noktada eğim sıfırdır: $f'(x) = 0$. Optimizasyonun
temel fikri budur --- kaybın türevinin sıfır olduğu parametreleri ararız.

\begin{center}
\begin{tabularx}{\textwidth}{@{}>{$}l<{$} >{$}l<{$} X@{}}
\toprule
\rowcolor{acikMavi}\color{anaMavi}\textbf{f(x)} & \color{anaMavi}\textbf{f'(x)} & \color{anaMavi}\textbf{Kural}\\
\midrule
x^n & n x^{n-1} & Kuvvet kuralı\\
\rowcolor{acikGri} e^x & e^x & Üstel\\
\ln x & 1/x & Logaritma\\
\rowcolor{acikGri} f(g(x)) & f'(g(x))\,g'(x) & Zincir kuralı (backprop'un temeli)\\
\bottomrule
\end{tabularx}
\end{center}

\gsub{Kısmi Türev ve Gradyan}

Makine öğrenmesinde kayıp, \emph{birçok} parametreye bağlıdır:
$L(w_1, w_2, \dots, w_n)$. Bir parametreye göre \emph{kısmi türev}
$\frac{\partial L}{\partial w_i}$, diğerlerini sabit tutup yalnızca $w_i$'yi
değiştirdiğimizde kaybın nasıl değiştiğini söyler. Tüm kısmi türevleri bir
vektörde toplarsak \emph{gradyanı} elde ederiz:
\[
\nabla L = \left[ \frac{\partial L}{\partial w_1},\;
\frac{\partial L}{\partial w_2},\; \dots,\;
\frac{\partial L}{\partial w_n} \right]^\top .
\]
Gradyan, kaybın \emph{en hızlı arttığı} yönü gösteren bir vektördür.
Dolayısıyla kaybı azaltmak için gradyanın \emph{tersi} yönde ilerleriz.

\gsub{Zincir Kuralı: Backprop'un Kalbi}

Zincir kuralı, \emph{bileşke fonksiyonların} türevini verir ve derin
öğrenmedeki geri yayılımın (backpropagation) temelidir. İç içe fonksiyonlarda
türevler \emph{çarpılarak} dışarıdan içeriye taşınır:
\[
\frac{d}{dx}f(g(x)) = f'(g(x))\cdot g'(x).
\]

\turetme{\textbf{Örnek:} $f(x) = (3x+1)^2$ fonksiyonunun türevi. Dış fonksiyon
$(\cdot)^2$, iç fonksiyon $3x+1$. Zincir kuralı:
\[
f'(x) = \underbrace{2(3x+1)}_{\text{dışın türevi}}\cdot
\underbrace{3}_{\text{için türevi}} = 6(3x+1).
\]
Daha uzun bir zincir, örneğin sigmoid içeren bir kayıp
$L = (\sigma(wx))^2$, aynı mantıkla katman katman çarpılır:
$\frac{dL}{dw} = 2\sigma(wx)\cdot\sigma'(wx)\cdot x$. Sinir ağları tam olarak
bu ``çarpım zincirini'' otomatikleştirir.}

\gsub{Kısmi Türev ve Gradyanın Somut Hesabı}

Kısmi türev kavramını sayısal bir örnekle pekiştirelim. Bir kaybın gradyanını
\emph{elle} hesaplamak, gradient descent'in ne yaptığını netleştirir.

\turetme{\textbf{İki parametreli kayıp:} $L(w_1, w_2) = (w_1 - 2)^2 + (w_2 + 1)^2$.
Bu, minimumu $(2, -1)$'de olan bir çanaktır. Kısmi türevler:
\[
\frac{\partial L}{\partial w_1} = 2(w_1 - 2), \qquad
\frac{\partial L}{\partial w_2} = 2(w_2 + 1).
\]
$(w_1, w_2) = (0, 0)$ noktasında gradyan:
\[
\nabla L = \big[\,2(0-2),\ 2(0+1)\,\big]^\top = [-4,\ 2]^\top.
\]
$\eta = 0.1$ ile bir gradient descent adımı:
\[
\vek{w} \leftarrow [0,0]^\top - 0.1\cdot[-4,2]^\top = [0.4,\ -0.2]^\top.
\]
Görüldüğü gibi $w_1$ minimuma ($2$) doğru arttı, $w_2$ minimuma ($-1$) doğru
azaldı --- tam istediğimiz yön.}

\begin{lstlisting}
import numpy as np

# L(w) = (w1 - 2)^2 + (w2 + 1)^2  kaybının gradyanı
def grad_L(w):
    return np.array([2 * (w[0] - 2), 2 * (w[1] + 1)])

w = np.array([0.0, 0.0])            # başlangıç noktası
lr = 0.1
for step in range(30):
    w = w - lr * grad_L(w)          # gradyanın tersine adım
print(w)                            # [2., -1.]'e yaklaşır (minimum)
\end{lstlisting}

\gsub{Gradient Descent (Gradyan İnişi)}

Gradient descent, kaybı minimize eden temel algoritmadır. Fikir basittir:
parametreleri, gradyanın ters yönünde küçük adımlarla güncelle. Güncelleme
kuralı:
\[
\vek{w} \leftarrow \vek{w} - \eta \, \nabla L(\vek{w}),
\]
burada $\eta > 0$ \emph{öğrenme oranıdır} (learning rate) --- adım büyüklüğünü
kontrol eder. Bunu, sisli bir dağda aşağı inmeye benzet: her adımda ayağının
altındaki en dik iniş yönünü bul ve o yöne bir adım at.

\turetme{$f(w) = w^2$ fonksiyonunu ele alalım; minimumu $w=0$'dadır. Türevi
$f'(w) = 2w$. $w_0 = 5$ ve $\eta = 0.1$ ile başlayalım:
\[
w_1 = 5 - 0.1 \cdot (2\cdot 5) = 5 - 1 = 4,\quad
w_2 = 4 - 0.1 \cdot 8 = 3.2,\quad
w_3 = 3.2 - 0.1 \cdot 6.4 = 2.56, \dots
\]
Her adımda $w$, minimuma ($0$) biraz daha yaklaşır. Adımlar minimuma
yaklaştıkça (gradyan küçüldükçe) küçülür.}

\begin{lstlisting}
import numpy as np

def gradient_descent(grad_fn, start, learning_rate=0.1, n_iters=50):
    """Genel amaçlı gradyan inişi.
    grad_fn: bir noktadaki gradyanı döndüren fonksiyon.
    """
    w = start
    history = [w]
    for _ in range(n_iters):
        grad = grad_fn(w)               # gradyanı hesapla
        w = w - learning_rate * grad    # gradyanın tersine adım at
        history.append(w)
    return w, history

# f(w) = w^2'yi minimize et; gradyanı f'(w) = 2w
grad = lambda w: 2 * w
w_min, hist = gradient_descent(grad, start=5.0)
print(w_min)                            # 0.0'a yaklaşır
\end{lstlisting}

\uyari{\textbf{Öğrenme oranı $\eta$ kritiktir.} Çok küçükse öğrenme çok yavaş
olur (binlerce adım). Çok büyükse adımlar minimumu aşar ve algoritma
\emph{ıraksar} (loss patlar, NaN olur). Pratikte $\eta$ genelde $0.001$--$0.1$
arasında denenir ve kayıp eğrisi izlenerek ayarlanır.}

\gsub{Gradient Descent Türleri}

\begin{center}
\begin{tabularx}{\textwidth}{@{}>{\bfseries\color{anaMavi}}l X X@{}}
\toprule
\rowcolor{acikMavi}\color{anaMavi}\textbf{Tür} & \color{anaMavi}\textbf{Her adımda kaç örnek?} & \color{anaMavi}\textbf{Özellik}\\
\midrule
Batch GD & Tüm veri ($m$) & Kararlı ama yavaş, çok bellek.\\
\rowcolor{acikGri} Stochastic (SGD) & Tek örnek (1) & Hızlı, gürültülü, kaçış yapabilir.\\
Mini-batch & Küçük grup (32--256) & En yaygın; denge; GPU dostu.\\
\bottomrule
\end{tabularx}
\end{center}

\gsub{Konvekslik: Neden Bazı Problemler Kolaydır?}

Bir fonksiyon \emph{konveksse} (dışbükey --- çanak şeklinde) tek bir global
minimumu vardır ve gradient descent oraya ulaşacağı garanti edilir. Lineer ve
lojistik regresyonun kayıpları konvekstir --- bu yüzden ``kolay'' problemlerdir.
Sinir ağlarının kaybı ise konveks \emph{değildir}: birçok yerel minimum ve eyer
noktası vardır; yine de gradient descent pratikte yeterince iyi çözümler bulur.

\bilgi{\textbf{Otomatik türev (autograd):} Modern kütüphaneler (PyTorch,
TensorFlow, JAX) gradyanları elle hesaplamana gerek bırakmaz; hesaplama
çizgesini takip edip zincir kuralını otomatik uygular. Yine de bu rehberde
gradyanları elle türeteceğiz --- çünkü ``sihrin'' altında ne olduğunu anlamak,
hata ayıklamak ve yeni modeller tasarlamak için şarttır.}

% =====================================================================
\gsec{Olasılık ve İstatistik}
% =====================================================================

Makine öğrenmesi belirsizlikle uğraşır: veri gürültülüdür, tahminler asla
kesin değildir. Olasılık teorisi, bu belirsizliği matematiksel olarak modeller.
Birçok ML algoritması (Naive Bayes, lojistik regresyon, tüm üretici modeller)
doğrudan olasılık üzerine kuruludur.

\gsub{Temel Kavramlar}

Bir \emph{rastgele değişken} $X$, sonucu belirsiz bir niceliktir. Olasılık
dağılımı, $X$'in hangi değerleri hangi olasılıkla aldığını tarif eder.
Ayrık değişkenler için \emph{olasılık kütle fonksiyonu} (PMF) $P(X=x)$,
sürekli için \emph{olasılık yoğunluk fonksiyonu} (PDF) $p(x)$ kullanılır.
Temel büyüklükler:
\[
\mathbb{E}[X] = \sum_x x\,P(X=x) \quad(\text{beklenen değer, ortalama}),
\qquad
\mathrm{Var}(X) = \mathbb{E}[(X - \mathbb{E}[X])^2].
\]
Beklenen değer dağılımın ``merkezini'' (ortalama), varyans ise ``yayılımını''
(dağınıklık) ölçer. Standart sapma $\sigma = \sqrt{\mathrm{Var}(X)}$'dır.

\turetme{\textbf{Örnek --- adil zar:} $X \in \{1,\dots,6\}$, her yüz olasılığı
$\frac{1}{6}$. Beklenen değer:
\[
\mathbb{E}[X] = \tfrac{1}{6}(1+2+3+4+5+6) = \tfrac{21}{6} = 3.5.
\]
Varyans $\mathbb{E}[X^2] - (\mathbb{E}[X])^2$ ile:
$\mathbb{E}[X^2]=\frac{1}{6}(1+4+9+16+25+36)=\frac{91}{6}\approx 15.17$, dolayısıyla
$\mathrm{Var}(X)=15.17 - 3.5^2 = 2.92$.}

\gsub{Önemli Olasılık Dağılımları}

ML'de sürekli karşılaşılan birkaç dağılımı tanımak, hangi modelin hangi
varsayıma dayandığını anlamak için gereklidir.

\begin{center}
\begin{tabularx}{\textwidth}{@{}>{\bfseries\color{anaMavi}}l >{$}l<{$} X@{}}
\toprule
\rowcolor{acikMavi}\color{anaMavi}\textbf{Dağılım} & \color{anaMavi}\textbf{Kütle/Yoğunluk} & \color{anaMavi}\textbf{Ne modeller?}\\
\midrule
Bernoulli & p^y(1-p)^{1-y} & Tek ikili sonuç (0/1). Lojistik regresyon.\\
\rowcolor{acikGri} Binom & \binom{n}{k}p^k(1-p)^{n-k} & $n$ denemede başarı sayısı.\\
Kategorik & \prod_c p_c^{[y=c]} & Çok sınıflı etiket. Softmax çıktısı.\\
\rowcolor{acikGri} Gaussian & \frac{1}{\sqrt{2\pi\sigma^2}}e^{-\frac{(x-\mu)^2}{2\sigma^2}} & Sürekli, gürültü. En küçük kareler.\\
\bottomrule
\end{tabularx}
\end{center}

\emph{Normal (Gaussian) dağılım} özellikle önemlidir: iki parametresi vardır ---
ortalama $\mu$ (merkez) ve standart sapma $\sigma$ (genişlik). Doğadaki birçok
gürültü normal dağıldığından, birazdan göreceğimiz gibi \emph{en küçük kareler
kaybı} bu varsayımdan doğar.

\begin{lstlisting}
import numpy as np

rng = np.random.default_rng(42)          # tekrarlanabilirlik için tohum

bernoulli = rng.binomial(1, 0.3, size=1000)   # p=0.3 ile 0/1 örnekler
gaussian = rng.normal(loc=0.0, scale=1.0, size=1000)  # mu=0, sigma=1

print(bernoulli.mean())                  # ~0.3 (ampirik p)
print(gaussian.mean(), gaussian.std())   # ~0.0, ~1.0
\end{lstlisting}

\gsub{Bağımsızlık ve Olasılık Zincir Kuralı}

İki olay \emph{bağımsızsa}, birlikte olma olasılıkları çarpımdır:
$P(A \cap B) = P(A)\,P(B)$. Bu basit kural, Naive Bayes'in ``saf''
varsayımının (özellikler sınıf verildiğinde bağımsız) ve MLE'de olabilirliğin
neden bir \emph{çarpım} olduğunun temelidir. Genel olarak (bağımlı olsalar bile)
zincir kuralı geçerlidir: $P(A \cap B) = P(A)\,P(B \mid A)$.

\gsub{Koşullu Olasılık ve Bayes Teoremi}

Koşullu olasılık $P(A \mid B)$, ``$B$ gerçekleştiği bilindiğinde $A$'nın
olasılığı''dır. Buradan makine öğrenmesinin en önemli formüllerinden biri olan
\emph{Bayes teoremi} çıkar:
\[
P(A \mid B) = \frac{P(B \mid A)\,P(A)}{P(B)}.
\]
Sınıflandırma diliyle: bir örneğin özellikleri $\vek{x}$ verildiğinde, sınıf
$y$ olma olasılığı ($\emph{posterior}$):
\[
\underbrace{P(y \mid \vek{x})}_{\text{sonsal}} =
\frac{\overbrace{P(\vek{x} \mid y)}^{\text{olabilirlik}}\;
\overbrace{P(y)}^{\text{önsel}}}{\underbrace{P(\vek{x})}_{\text{kanıt}}}.
\]
Bu formül, Naive Bayes sınıflandırıcısının (Kısım 4) doğrudan temelidir:
her sınıf için sağ tarafı hesaplar ve en yüksek olasılıklı sınıfı seçeriz.

\turetme{\textbf{Örnek --- tıbbi test:} Bir hastalık nüfusun \%1'inde var
($P(H)=0.01$). Test hastalığı \%99 yakalıyor ($P(+\mid H)=0.99$) ama sağlıklıda
da \%5 yanlış alarm veriyor ($P(+\mid \bar H)=0.05$). Testin pozitif çıktı;
gerçekten hasta olma olasılığın nedir?
\[
P(H\mid +) = \frac{P(+\mid H)P(H)}{P(+\mid H)P(H)+P(+\mid \bar H)P(\bar H)}
= \frac{0.99 \cdot 0.01}{0.99\cdot 0.01 + 0.05 \cdot 0.99}
\approx 0.167.
\]
Yani pozitif teste rağmen gerçekten hasta olma olasılığın yalnızca \%17!
Nadir olayların (düşük önsel) sezgiye aykırı sonuçları buradan gelir.}

\gsub{Maksimum Olabilirlik Tahmini (MLE)}

Maksimum olabilirlik (Maximum Likelihood Estimation), bir modelin
parametrelerini veriden tahmin etmenin temel yöntemidir ve birçok kayıp
fonksiyonunun \emph{nereden geldiğini} açıklar. Fikir: gözlemlediğimiz veriyi
\emph{en olası} kılan parametreleri seç.

Bağımsız gözlemler $x_1, \dots, x_m$ için olabilirlik, çarpımdır:
\[
\mathcal{L}(\theta) = \prod_{i=1}^{m} p(x_i \mid \theta).
\]
Çarpımlarla çalışmak zor olduğundan (ve sayısal olarak taşabildiğinden)
\emph{log-olabilirliği} maksimize ederiz (log tekdüze arttığından maksimum
noktası aynıdır):
\[
\ell(\theta) = \log \mathcal{L}(\theta) = \sum_{i=1}^{m} \log p(x_i \mid \theta).
\]

\turetme{\textbf{Basit örnek --- hileli para:} Bir parayı 10 kez attık ve 7 kez
yazı geldi. Yazı olasılığı $p$'yi MLE ile tahmin edelim. Olabilirlik
(Bernoulli): $\mathcal{L}(p) = p^7 (1-p)^3$. Log-olabilirlik:
\[
\ell(p) = 7\ln p + 3\ln(1-p).
\]
Maksimum için türevi sıfırlarız:
\[
\ell'(p) = \frac{7}{p} - \frac{3}{1-p} = 0
\;\Longrightarrow\; 7(1-p) = 3p
\;\Longrightarrow\; p = \frac{7}{10}.
\]
Yani MLE tahmini $\hat{p}=0.7$ --- tam da ampirik frekans (yazı sayısı / toplam).
Sezgimizle örtüşür; MLE'nin gücü, sezgisel olmayan modellerde de aynı ilkeyi
uygulayabilmesidir.}

\turetme{\textbf{MLE $\Rightarrow$ En Küçük Kareler:} Regresyonda hataların
normal dağıldığını, $y_i = \vek{w}^\top\vek{x}_i + \varepsilon_i$ ve
$\varepsilon_i \sim \mathcal{N}(0,\sigma^2)$, varsayalım. Bir örneğin
olabilirliği:
\[
p(y_i \mid \vek{x}_i) = \frac{1}{\sqrt{2\pi\sigma^2}}
\exp\!\left(-\frac{(y_i - \vek{w}^\top\vek{x}_i)^2}{2\sigma^2}\right).
\]
Log-olabilirlik toplandığında, $\vek{w}$'ye bağlı olmayan sabitler atılınca:
\[
\ell(\vek{w}) = -\frac{1}{2\sigma^2}\sum_{i=1}^{m}(y_i - \vek{w}^\top\vek{x}_i)^2 + \text{sabit}.
\]
$\ell$'i maksimize etmek, $\sum (y_i - \vek{w}^\top\vek{x}_i)^2$ toplamını
\emph{minimize} etmekle aynıdır --- yani \textbf{ortalama kare hatası}!
Demek ki en küçük kareler, Gauss gürültü varsayımı altında MLE'dir.}

\bilgi{Bu türetme çok önemli bir dersi öğretir: makine öğrenmesindeki kayıp
fonksiyonları keyfi değildir. Ortalama kare hatası (MSE) normal dağılımdan,
çapraz entropi (cross-entropy) Bernoulli/kategorik dağılımdan MLE yoluyla
doğar. Kaybı seçmek, aslında verinin dağılımı hakkında bir varsayım yapmaktır.}

% =====================================================================
\gsec{Bilgi Teorisi}
% =====================================================================

Bilgi teorisi, ``belirsizliği'' ve ``bilgiyi'' niceliksel ölçer. Sınıflandırma
kayıplarının (cross-entropy) ve karar ağaçlarının bölme kriterlerinin
(information gain) temelini oluşturur.

\gsub{Entropi}

Entropi, bir olasılık dağılımının \emph{belirsizliğini} ölçer. Bir olayın
taşıdığı ``şaşırtıcılık'' $-\log p$ ile tanımlanır (nadir olaylar daha
şaşırtıcı). Entropi, bu şaşırtıcılığın beklenen değeridir:
\[
H(P) = -\sum_{i} p_i \log p_i .
\]
Dağılım ne kadar düzgünse (belirsizse) entropi o kadar yüksektir. Adil bir
paranın entropisi maksimumdur ($\log 2 = 1$ bit); hep yazı gelen hileli bir
paranın entropisi sıfırdır (belirsizlik yok).

\begin{lstlisting}
import numpy as np

def entropy(p):
    """Bir olasılık dağılımının Shannon entropisi (bit cinsinden)."""
    p = np.asarray(p)
    p = p[p > 0]                    # log(0)'dan kaçın
    return -np.sum(p * np.log2(p))

print(entropy([0.5, 0.5]))          # 1.0 bit  (maksimum belirsizlik)
print(entropy([0.9, 0.1]))          # 0.469 bit
print(entropy([1.0, 0.0]))          # 0.0 bit  (belirsizlik yok)
\end{lstlisting}

\gsub{Çapraz Entropi ve KL Iraksaması}

\emph{Çapraz entropi}, gerçek dağılım $P$ yerine tahmin edilen dağılım $Q$
kullanıldığında gereken ortalama şaşırtıcılıktır:
\[
H(P, Q) = -\sum_i p_i \log q_i .
\]
Sınıflandırmada $P$ gerçek etiket (one-hot), $Q$ modelin tahmin ettiği
olasılıklardır. Çapraz entropiyi minimize etmek, modelin tahminlerini gerçek
etiketlere yaklaştırır --- bu yüzden sınıflandırmanın standart kaybıdır.

İkili sınıflandırmada ($y \in \{0,1\}$) çapraz entropi, tek bir örnek için
şu tanıdık biçimi alır (\emph{binary cross-entropy}):
\[
L = -\big[\,y\log\hat{y} + (1-y)\log(1-\hat{y})\,\big],
\]
burada $\hat{y}$ modelin ``sınıf $1$'' olma tahminidir. Kısım 3'te bunun
Bernoulli MLE'sinden doğduğunu göreceğiz.

\turetme{\textbf{Örnek --- çapraz entropi hesabı:} Gerçek etiket 3 sınıflı,
one-hot $P=(0,1,0)$ (yani doğru sınıf ikinci). Model iki farklı tahmin
üretsin:
\begin{itemize}\itemsep1pt
\item İyi model $Q=(0.1,\,0.8,\,0.1)$: \ $L = -\ln(0.8) \approx 0.22$.
\item Kötü model $Q=(0.4,\,0.3,\,0.3)$: \ $L = -\ln(0.3) \approx 1.20$.
\end{itemize}
Yalnızca doğru sınıfa atanan olasılığın logaritması sayılır (diğer terimler
$p_i=0$ ile çarpılıp yok olur). Model doğru sınıfa ne kadar yüksek olasılık
verirse kayıp o kadar küçük olur --- doğruya $1$ verirse kayıp $0$.}

\begin{lstlisting}
import numpy as np

def cross_entropy(p_true, q_pred):
    """One-hot gerçek dağılım ile tahmin arası çapraz entropi (doğal log)."""
    q_pred = np.clip(q_pred, 1e-12, 1.0)     # log(0)'ı önle
    return -np.sum(p_true * np.log(q_pred))

p = np.array([0, 1, 0])                      # gerçek sınıf: ikinci
print(cross_entropy(p, [0.1, 0.8, 0.1]))     # ~0.22 (iyi tahmin)
print(cross_entropy(p, [0.4, 0.3, 0.3]))     # ~1.20 (kötü tahmin)
\end{lstlisting}

\emph{Kullback-Leibler (KL) ıraksaması}, iki dağılım arasındaki ``uzaklığı''
ölçer:
\[
D_{\mathrm{KL}}(P \parallel Q) = \sum_i p_i \log \frac{p_i}{q_i}
= H(P,Q) - H(P).
\]
KL her zaman $\ge 0$'dır ve yalnızca $P=Q$ iken sıfırdır. Gerçek dağılım $P$
sabit olduğunda $H(P)$ de sabittir; dolayısıyla \emph{çapraz entropiyi
minimize etmek, KL ıraksamasını minimize etmekle aynıdır}.

\begin{center}
\begin{tabularx}{\textwidth}{@{}>{\bfseries\color{anaMavi}}l X X@{}}
\toprule
\rowcolor{acikMavi}\color{anaMavi}\textbf{Kavram} & \color{anaMavi}\textbf{Ölçtüğü} & \color{anaMavi}\textbf{ML'de kullanımı}\\
\midrule
Entropi & Bir dağılımın belirsizliği & Karar ağacı bölme (information gain)\\
\rowcolor{acikGri} Çapraz entropi & $Q$ ile $P$'yi kodlama maliyeti & Sınıflandırma kaybı\\
KL ıraksaması & İki dağılım arası fark & VAE, dağılım eşleştirme\\
\bottomrule
\end{tabularx}
\end{center}

\notk[Bağlantı]{Kısım 3'te lojistik regresyonun kaybını türettiğimizde,
Bernoulli dağılımı için MLE'nin tam olarak \emph{ikili çapraz entropiye}
eşit olduğunu göreceğiz. Böylece olasılık (MLE), bilgi teorisi (cross-entropy)
ve optimizasyon (gradient descent) tek bir noktada birleşecek. Bu üç temelin
nasıl kaynaştığını görmek, makine öğrenmesini gerçekten ``anlamak'' demektir.}

% #####################################################################
\gpart{2}{Öğrenme Teorisi}
% #####################################################################

% =====================================================================
\gsec{Öğrenme Problemi ve Kayıp Fonksiyonları}
% =====================================================================

Makine öğrenmesini kurmadan önce ``öğrenmenin'' ne demek olduğunu matematiksel
olarak tanımlamalıyız. Denetimli öğrenmede elimizde girdi-çıktı çiftleri
$\{(\vek{x}_i, y_i)\}_{i=1}^{m}$ vardır ve amacımız, yeni bir $\vek{x}$ için
$y$'yi tahmin eden bir $f_{\vek{\theta}}(\vek{x})$ fonksiyonu (model) bulmaktır.

\gsub{Temel Bileşenler}

\begin{center}
\begin{tabularx}{\textwidth}{@{}>{\bfseries\color{anaMavi}}l X@{}}
\toprule
\rowcolor{acikMavi}\color{anaMavi}\textbf{Bileşen} & \color{anaMavi}\textbf{Tanım}\\
\midrule
Model $f_{\vek{\theta}}$ & Girdiden çıktıya eşleyen, $\vek{\theta}$ parametreli fonksiyon.\\
\rowcolor{acikGri} Kayıp $L$ & Tek bir tahminin ne kadar yanlış olduğu.\\
Maliyet $J$ & Tüm veri üzerinde ortalama kayıp: $J(\vek{\theta})=\frac1m\sum_i L(f_{\vek\theta}(\vek x_i), y_i)$.\\
\rowcolor{acikGri} Optimizasyon & $J$'yi minimize eden $\vek{\theta}$'yı bulma (öğrenme budur).\\
\bottomrule
\end{tabularx}
\end{center}

\gsub{Yaygın Kayıp Fonksiyonları}

Kayıp fonksiyonu, ``iyi tahmin'' tanımını belirler ve probleme göre seçilir.

\begin{center}
\begin{tabularx}{\textwidth}{@{}>{\bfseries\color{anaMavi}}l >{$}l<{$} X@{}}
\toprule
\rowcolor{acikMavi}\color{anaMavi}\textbf{Kayıp} & \color{anaMavi}\textbf{Formül} & \color{anaMavi}\textbf{Kullanım}\\
\midrule
Kare hata (MSE) & (y - \hat y)^2 & Regresyon\\
\rowcolor{acikGri} Mutlak hata (MAE) & |y - \hat y| & Regresyon (aykırıya dayanıklı)\\
İkili çapraz entropi & -[y\log\hat y + (1{-}y)\log(1{-}\hat y)] & İkili sınıflandırma\\
\rowcolor{acikGri} Hinge & \max(0, 1 - y\hat y) & SVM\\
\bottomrule
\end{tabularx}
\end{center}

\gsub{Ampirik Risk Minimizasyonu}

Gerçekte veriyi üreten dağılımı bilmeyiz; elimizde yalnızca örneklem vardır.
Bu yüzden ``gerçek riski'' (tüm dağılım üzerinde beklenen kayıp) doğrudan
minimize edemeyiz; onun yerine \emph{ampirik riski} (eğitim verisi üzerindeki
ortalama kayıp) minimize ederiz:
\[
\hat{\vek{\theta}} = \arg\min_{\vek{\theta}} \; \frac{1}{m}
\sum_{i=1}^{m} L\big(f_{\vek{\theta}}(\vek{x}_i), y_i\big).
\]
Buradaki tehlike şudur: eğitim verisinde mükemmel olmak, \emph{yeni} veride
iyi olmayı garanti etmez. Bu, aşırı öğrenme problemidir.

% =====================================================================
\gsec{Aşırı Öğrenme, Bias-Variance ve Düzenlileştirme}
% =====================================================================

\gsub{Aşırı ve Eksik Öğrenme}

Bir modelin başarısı, eğitim verisini ezberlemesiyle değil, \emph{yeni} veriye
\emph{genelleme} yapmasıyla ölçülür. İki başarısızlık modu vardır:

\begin{center}
\begin{tabularx}{\textwidth}{@{}>{\bfseries\color{anaMavi}}l X X@{}}
\toprule
\rowcolor{acikMavi}\color{anaMavi}\textbf{Durum} & \color{anaMavi}\textbf{Belirti} & \color{anaMavi}\textbf{Neden / çözüm}\\
\midrule
Eksik öğrenme & Eğitim VE test kötü & Model çok basit; daha güçlü model / özellik.\\
\rowcolor{acikGri} İyi denge & Her ikisi de iyi & İstenen nokta.\\
Aşırı öğrenme & Eğitim iyi, test kötü & Model çok karmaşık; düzenlileştirme / daha çok veri.\\
\bottomrule
\end{tabularx}
\end{center}

\gsub{Bias-Variance Ayrışması}

Bir modelin beklenen test hatası matematiksel olarak üç parçaya ayrılır:
\[
\mathbb{E}[(y - \hat f(\vek{x}))^2] =
\underbrace{\mathrm{Bias}[\hat f]^2}_{\text{yanlılık}} +
\underbrace{\mathrm{Var}[\hat f]}_{\text{varyans}} +
\underbrace{\sigma^2}_{\text{indirgenemez gürültü}}.
\]
\textbf{Bias (yanlılık):} modelin sistematik hatası --- çok basit bir model
gerçeği kaçırır (eksik öğrenme). \textbf{Variance (varyans):} modelin eğitim
verisindeki küçük değişimlere aşırı duyarlılığı --- çok karmaşık model gürültüyü
öğrenir (aşırı öğrenme). İkisini aynı anda düşürmek genelde mümkün değildir;
bu \emph{bias-variance ödünleşmesidir}.

\bilgi{Model karmaşıklığı arttıkça bias düşer ama variance artar. Toplam hata
önce düşer, bir minimumdan sonra tekrar yükselir. Amaç, bu \emph{tatlı noktayı}
bulmaktır --- ne fazla basit ne fazla karmaşık. Düzenlileştirme ve çapraz
doğrulama bunun için araçlardır.}

\gsub{Düzenlileştirme (Regularization)}

Düzenlileştirme, aşırı öğrenmeyi önlemek için maliyet fonksiyonuna model
karmaşıklığını cezalandıran bir terim ekler. Böylece model, veriye uymakla
\emph{basit kalmak} arasında denge kurar:
\[
J(\vek{w}) = \underbrace{\frac{1}{m}\sum_{i=1}^{m} L(f_{\vek w}(\vek x_i), y_i)}_{\text{veriye uyum}}
+ \lambda \underbrace{R(\vek{w})}_{\text{karmaşıklık cezası}}.
\]
Burada $\lambda \ge 0$ ceza gücünü ayarlar. İki temel ceza:
\[
\text{Ridge (L2):}\; R(\vek w) = \|\vek w\|_2^2 = \sum_j w_j^2, \qquad
\text{Lasso (L1):}\; R(\vek w) = \|\vek w\|_1 = \sum_j |w_j|.
\]
L2 ağırlıkları küçültür ama sıfırlamaz; L1 bazı ağırlıkları \emph{tam sıfır}
yapar, böylece otomatik özellik seçimi sağlar. Bu farkın nedenini Kısım 3'te
göreceğiz.

\uyari{$\lambda$ bir hiperparametredir --- veriden öğrenilmez, çapraz
doğrulamayla seçilir. $\lambda$ çok büyükse model aşırı basitleşir (eksik
öğrenme, yüksek bias); çok küçükse düzenlileştirme etkisiz kalır (aşırı öğrenme,
yüksek variance). Doğru $\lambda$'yı bir doğrulama kümesinde aranmalıdır.}

% #####################################################################
\gpart{3}{Regresyon}
% #####################################################################

% =====================================================================
\gsec{Lineer Regresyon}
% =====================================================================

Lineer regresyon, makine öğrenmesinin ``merhaba dünya''sıdır: basit ama
üzerine çok şey inşa edilen temel bir modeldir. Çıktının, girdilerin
\emph{ağırlıklı toplamı} olduğunu varsayar.

\gsub{Model}

$n$ özellikli bir örnek $\vek{x}=(x_1,\dots,x_n)$ için tahmin:
\[
\hat y = w_0 + w_1 x_1 + \dots + w_n x_n = \vek{w}^\top \vek{x},
\]
burada $w_0$ bias (kesişim) terimidir; $\vek{x}$'e sabit bir $1$ eklenerek
(bias trick) matris çarpımına dahil edilir. Tüm veri için:
$\hat{\vek{y}} = \vek{X}\vek{w}$, $\vek{X}\in\R^{m\times(n+1)}$.

\gsub{Maliyet Fonksiyonu}

Kısım 1'de MLE'den türettiğimiz gibi, kayıp ortalama kare hatasıdır:
\[
J(\vek{w}) = \frac{1}{2m}\sum_{i=1}^{m}\big(\vek{w}^\top\vek{x}_i - y_i\big)^2
= \frac{1}{2m}\|\vek{X}\vek{w} - \vek{y}\|_2^2 .
\]
$\frac{1}{2}$ çarpanı yalnızca türevi sadeleştirmek içindir (sonuçları
değiştirmez). Bu $J$ konvekstir, dolayısıyla tek bir global minimumu vardır.

\gsub{Çözüm 1: Normal Denklem (Kapalı Form)}

Konveks $J$'nin minimumunu gradyanını sıfıra eşitleyerek \emph{analitik olarak}
bulabiliriz.

\turetme{Maliyeti matris biçiminde açalım:
\[
J(\vek w) = \frac{1}{2m}(\vek{X}\vek{w}-\vek{y})^\top(\vek{X}\vek{w}-\vek{y}).
\]
$\vek w$'ye göre gradyan (matris kalkülüsü kuralları ile):
\[
\nabla_{\vek w} J = \frac{1}{m}\vek{X}^\top(\vek{X}\vek{w}-\vek{y}).
\]
Sıfıra eşitleyip çözersek \emph{normal denklemi} elde ederiz:
\[
\vek{X}^\top\vek{X}\vek{w} = \vek{X}^\top\vek{y}
\;\;\Longrightarrow\;\;
\boxed{\;\vek{w} = (\vek{X}^\top\vek{X})^{-1}\vek{X}^\top\vek{y}\;}.
\]}

\begin{lstlisting}
import numpy as np

class LinearRegressionClosedForm:
    def fit(self, X, y):
        # bir'lerden oluşan bias sütunu ekle -> X (m, n+1) olur
        X_b = np.c_[np.ones((X.shape[0], 1)), X]
        # Normal denklem: w = (X^T X)^{-1} X^T y
        self.w = np.linalg.inv(X_b.T @ X_b) @ X_b.T @ y
        return self

    def predict(self, X):
        X_b = np.c_[np.ones((X.shape[0], 1)), X]
        return X_b @ self.w

# Örnek
X = np.array([[1.0], [2.0], [3.0], [4.0]])
y = np.array([2.0, 4.0, 6.0, 8.0])          # tam olarak y = 2x
model = LinearRegressionClosedForm().fit(X, y)
print(model.w)                               # ~[0, 2]  (bias 0, eğim 2)
\end{lstlisting}

\uyari{Normal denklem $(\vek{X}^\top\vek{X})^{-1}$ matris tersi gerektirir;
bu $\mathcal{O}(n^3)$'tür ve özellik sayısı $n$ çok büyükse (binlerce) pahalıdır.
Ayrıca $\vek{X}^\top\vek{X}$ tekil (tersinemez) olabilir. Bu durumlarda
gradient descent tercih edilir.}

\gsub{Çözüm 2: Gradient Descent}

Büyük veri için normal denklem yerine gradyan inişini kullanırız. Yukarıda
türettiğimiz gradyanı doğrudan kullanan güncelleme:
\[
\vek{w} \leftarrow \vek{w} - \eta \cdot \frac{1}{m}\vek{X}^\top(\vek{X}\vek{w}-\vek{y}).
\]

\begin{lstlisting}
class LinearRegressionGD:
    def __init__(self, lr=0.01, n_iters=1000):
        self.lr = lr
        self.n_iters = n_iters

    def fit(self, X, y):
        m, n = X.shape
        self.w = np.zeros(n)
        self.b = 0.0
        for _ in range(self.n_iters):
            y_pred = X @ self.w + self.b            # tahminler (m,)
            error = y_pred - y                      # artıklar (m,)
            # MSE'nin ağırlıklara ve bias'a göre gradyanları
            dw = (1 / m) * (X.T @ error)            # (n,)
            db = (1 / m) * np.sum(error)            # skaler
            self.w -= self.lr * dw                  # gradyan adımı
            self.b -= self.lr * db
        return self

    def predict(self, X):
        return X @ self.w + self.b
\end{lstlisting}

\gsub{scikit-learn ile}

\begin{lstlisting}
from sklearn.linear_model import LinearRegression
from sklearn.metrics import mean_squared_error, r2_score

model = LinearRegression()
model.fit(X_train, y_train)
y_pred = model.predict(X_test)

print(model.coef_, model.intercept_)        # öğrenilen ağırlıklar ve bias
print(r2_score(y_test, y_pred))             # uyum iyiliği
\end{lstlisting}

% =====================================================================
\gsec{Ridge, Lasso ve Lojistik Regresyon}
% =====================================================================

\gsub{Ridge Regresyonu (L2)}

Ridge, maliyete $L2$ cezası ekler. İlginç olan, kapalı formunun hâlâ mevcut
olması ve tekillik problemini de çözmesidir:
\[
J(\vek w) = \frac{1}{2m}\|\vek{X}\vek{w}-\vek{y}\|_2^2 + \frac{\lambda}{2}\|\vek{w}\|_2^2
\;\;\Longrightarrow\;\;
\vek{w} = (\vek{X}^\top\vek{X} + \lambda\vek{I})^{-1}\vek{X}^\top\vek{y}.
\]
$\lambda\vek{I}$ terimi köşegeni güçlendirir, matrisi her zaman tersinir kılar
--- bu yüzden Ridge sayısal olarak da kararlıdır.

\gsub{Neden Lasso Sıfırlar, Ridge Sıfırlamaz?}

Bu, sıkça sorulan ve derin bir sorudur. Geometrik sezgi: düzenlileştirme,
ağırlıkları bir ``bütçe bölgesi'' içinde kalmaya zorlar. $L2$ bölgesi bir
\emph{daire} (yuvarlak), $L1$ bölgesi bir \emph{karedir} (köşeli). Maliyet
konturları bu bölgeye ilk değdiği noktada çözüm bulunur. $L1$'in köşeleri
eksenler üzerindedir --- kontur genelde bir köşeye değer ve o eksendeki ağırlık
tam sıfır olur. $L2$'nin yuvarlak sınırında böyle köşe yoktur, bu yüzden
sıfırlamaz yalnızca küçültür.

\begin{center}
\begin{tabularx}{\textwidth}{@{}>{\bfseries\color{anaMavi}}l X X@{}}
\toprule
\rowcolor{acikMavi}\color{anaMavi}\textbf{} & \color{anaMavi}\textbf{Ridge (L2)} & \color{anaMavi}\textbf{Lasso (L1)}\\
\midrule
Ceza & $\sum w_j^2$ & $\sum |w_j|$\\
\rowcolor{acikGri} Ağırlıklara etki & Küçültür (sıfırlamaz) & Bazılarını sıfırlar\\
Özellik seçimi & Hayır & Evet (seyrek çözüm)\\
\rowcolor{acikGri} Kapalı form & Var & Yok (koordinat inişi)\\
\bottomrule
\end{tabularx}
\end{center}

\begin{lstlisting}
from sklearn.linear_model import Ridge, Lasso, ElasticNet

ridge = Ridge(alpha=1.0).fit(X_train, y_train)
lasso = Lasso(alpha=0.1).fit(X_train, y_train)     # bazı coef_ 0 olur
elastic = ElasticNet(alpha=0.1, l1_ratio=0.5)      # L1 + L2 birleşik
print(lasso.coef_)                                  # sıfırları gözlemle
\end{lstlisting}

\gsub{Lojistik Regresyon: Sınıflandırmaya Köprü}

Adına ``regresyon'' dense de lojistik regresyon bir \emph{sınıflandırma}
algoritmasıdır. Lineer bir skoru, \emph{sigmoid} fonksiyonuyla $[0,1]$
aralığına --- bir olasılığa --- sıkıştırır:
\[
\sigma(z) = \frac{1}{1 + e^{-z}}, \qquad
\hat y = P(y=1\mid\vek{x}) = \sigma(\vek{w}^\top\vek{x}).
\]
Sigmoid'in güzel bir özelliği türevidir: $\sigma'(z) = \sigma(z)(1-\sigma(z))$
--- bu, gradyanı sadeleştirir.

\gsub{Kayıp: İkili Çapraz Entropi (MLE'den)}

\turetme{Etiket $y\in\{0,1\}$ Bernoulli dağılır: $P(y\mid\vek x)=\hat y^{\,y}(1-\hat y)^{1-y}$.
Log-olabilirlik, tek örnek için:
\[
\log P(y\mid\vek x) = y\log\hat y + (1-y)\log(1-\hat y).
\]
Bunun \emph{negatifini} minimize etmek, tüm veri için maliyeti verir ---
\textbf{ikili çapraz entropi}:
\[
J(\vek w) = -\frac{1}{m}\sum_{i=1}^{m}\big[y_i\log\hat y_i + (1-y_i)\log(1-\hat y_i)\big].
\]
Gradyanı türetince (sigmoid türevi sayesinde) çarpıcı biçimde sadeleşir ---
lineer regresyonunkiyle \emph{aynı} biçimi alır:
\[
\nabla_{\vek w} J = \frac{1}{m}\vek{X}^\top(\hat{\vek y} - \vek{y}).
\]}

\begin{lstlisting}
import numpy as np

class LogisticRegressionScratch:
    def __init__(self, lr=0.1, n_iters=1000):
        self.lr, self.n_iters = lr, n_iters

    @staticmethod
    def _sigmoid(z):
        return 1.0 / (1.0 + np.exp(-z))

    def fit(self, X, y):
        m, n = X.shape
        self.w = np.zeros(n)
        self.b = 0.0
        for _ in range(self.n_iters):
            z = X @ self.w + self.b
            y_hat = self._sigmoid(z)                # tahmin edilen olasılıklar
            # Lineer regresyonla aynı gradyan biçimi (MLE türetmesiyle)
            dw = (1 / m) * (X.T @ (y_hat - y))
            db = (1 / m) * np.sum(y_hat - y)
            self.w -= self.lr * dw
            self.b -= self.lr * db
        return self

    def predict_proba(self, X):
        return self._sigmoid(X @ self.w + self.b)

    def predict(self, X, threshold=0.5):
        return (self.predict_proba(X) >= threshold).astype(int)
\end{lstlisting}

\bilgi{Lineer ve lojistik regresyonun gradyanının \emph{aynı biçimde}
($\vek{X}^\top(\hat{\vek y}-\vek y)$) çıkması tesadüf değildir: her ikisi de
``genelleştirilmiş lineer modeller'' ailesindendir. Bu ortak yapı, aynı
optimizasyon kodunu farklı modellerde yeniden kullanmanı sağlar ---
ve derin öğrenmedeki katman gradyanlarının da temelini oluşturur.}

% #####################################################################
\gpart{4}{Sınıflandırma Algoritmaları}
% #####################################################################

% =====================================================================
\gsec{Naive Bayes ve K-En Yakın Komşu}
% =====================================================================

\gsub{Naive Bayes}

Naive Bayes, Kısım 1'deki Bayes teoreminin doğrudan uygulamasıdır. Bir örneğin
her sınıfa ait olma olasılığını hesaplar ve en yükseğini seçer:
\[
\hat y = \arg\max_{y} \; P(y \mid \vek{x}) = \arg\max_{y}\; P(\vek{x}\mid y)\,P(y).
\]
``Naive'' (saf) olmasının nedeni, güçlü bir varsayımdır: özelliklerin, sınıf
verildiğinde \emph{birbirinden bağımsız} olduğunu kabul eder. Bu varsayım
sayesinde ortak olasılık, tek tek olasılıkların çarpımına iner:
\[
P(\vek{x}\mid y) = \prod_{j=1}^{n} P(x_j \mid y).
\]
Varsayım gerçekçi olmasa da (özellikler genelde ilişkilidir), pratikte ---
özellikle metin sınıflandırmada (spam filtresi) --- şaşırtıcı derecede iyi
çalışır ve çok hızlıdır.

\turetme{Sayısal taşmayı önlemek için çarpım yerine \emph{log} kullanılır
(log, çarpımı toplama çevirir ve tekdüze olduğundan argmax'i değiştirmez):
\[
\hat y = \arg\max_{y}\;\Big[\log P(y) + \sum_{j=1}^{n}\log P(x_j\mid y)\Big].
\]
Sürekli özellikler için $P(x_j\mid y)$ genelde normal (Gaussian) dağılımla
modellenir; her sınıf-özellik çifti için ortalama ve varyans eğitim verisinden
tahmin edilir.}

\begin{lstlisting}
from sklearn.naive_bayes import GaussianNB, MultinomialNB

gnb = GaussianNB()                          # sürekli özellikler
gnb.fit(X_train, y_train)
print(gnb.predict(X_test))

# Metin (kelime sayıları) için MultinomialNB standarttır:
# mnb = MultinomialNB(alpha=1.0)  # alpha = Laplace düzeltmesi
\end{lstlisting}

\uyari{\textbf{Sıfır olasılık problemi:} Eğitimde hiç görülmemiş bir
özellik-sınıf kombinasyonu olasılığı $0$ yapar ve tüm çarpımı sıfırlar. Çözüm
\emph{Laplace düzeltmesidir} (smoothing): her sayıma küçük bir sabit ($\alpha$)
eklenir, böylece hiçbir olasılık tam sıfır olmaz.}

\gsub{K-En Yakın Komşu (KNN)}

KNN, ``benzer örnekler benzer etikete sahiptir'' sezgisine dayanan, eğitimsiz
(lazy) bir algoritmadır. Yeni bir örneği sınıflandırmak için: eğitim
kümesindeki \emph{en yakın} $k$ komşuyu bul, çoğunluğun etiketini ver.
``Yakınlık'' genelde Öklid uzaklığıyla ölçülür:
\[
d(\vek{x}, \vek{x}') = \|\vek{x}-\vek{x}'\|_2 = \sqrt{\sum_{j=1}^{n}(x_j - x'_j)^2}.
\]
KNN'nin ``eğitimi'' yoktur --- tüm işi tahmin anında yapar (veriyi saklar,
sorguda uzaklıkları hesaplar). Bu yüzden büyük veride tahmin yavaştır.

\begin{lstlisting}
import numpy as np
from collections import Counter

class KNNClassifier:
    def __init__(self, k=5):
        self.k = k

    def fit(self, X, y):
        self.X_train = X            # tembel: veriyi sadece ezberle
        self.y_train = y
        return self

    def predict(self, X):
        preds = []
        for x in X:
            # her eğitim noktasına Öklid uzaklığı
            distances = np.sqrt(np.sum((self.X_train - x) ** 2, axis=1))
            k_idx = np.argsort(distances)[:self.k]      # en yakın k indeks
            k_labels = self.y_train[k_idx]
            preds.append(Counter(k_labels).most_common(1)[0][0])  # çoğunluk
        return np.array(preds)
\end{lstlisting}

\bilgi{\textbf{$k$ seçimi bir ödünleşmedir:} Küçük $k$ (örneğin 1) gürültüye
duyarlıdır (yüksek variance); büyük $k$ sınırları aşırı yumuşatır (yüksek bias).
KNN \emph{ölçeğe çok duyarlıdır} --- büyük değerli özellikler uzaklığa hâkim
olur; bu yüzden KNN'den önce mutlaka özellikleri standartlaştır.}

% =====================================================================
\gsec{Destek Vektör Makineleri (SVM)}
% =====================================================================

SVM, iki sınıfı ayıran en iyi \emph{hiperdüzlemi} bulan güçlü bir
sınıflandırıcıdır. ``En iyi'' derken: iki sınıf arasındaki boşluğu (marjı)
\emph{maksimize} eden düzlemi kasteder. Geniş marj, daha iyi genelleme demektir.

\gsub{Maksimum Marj}

Bir hiperdüzlem $\vek{w}^\top\vek{x} + b = 0$ ile tanımlanır. Sınıflar
$y\in\{-1,+1\}$ olsun. Her örneğin düzleme uzaklığının en azını (marjı)
maksimize etmek istiyoruz. Bu, şu kısıtlı optimizasyona dönüşür:
\[
\min_{\vek{w}, b} \; \frac{1}{2}\|\vek{w}\|_2^2
\quad \text{öyle ki} \quad
y_i(\vek{w}^\top\vek{x}_i + b) \ge 1, \;\; \forall i.
\]
Marj $\frac{2}{\|\vek{w}\|}$ olduğundan, onu maksimize etmek $\|\vek{w}\|$'yi
minimize etmekle aynıdır. Kısıtı sağlayan, marj sınırında yatan örneklere
\emph{destek vektörleri} denir --- çözümü yalnızca onlar belirler.

\gsub{Yumuşak Marj (Soft Margin)}

Gerçek veri genelde tam ayrılabilir değildir. \emph{Yumuşak marj}, bir miktar
ihlale izin verir; bunu $\xi_i$ ``gevşeklik'' değişkenleri ve bir $C$
parametresiyle dengeler:
\[
\min_{\vek{w},b,\vek{\xi}} \; \frac{1}{2}\|\vek{w}\|^2 + C\sum_{i=1}^m \xi_i,
\quad y_i(\vek{w}^\top\vek{x}_i+b)\ge 1-\xi_i,\;\; \xi_i\ge 0.
\]
$C$ ödünleşmeyi kontrol eder: büyük $C$ ihlalleri sert cezalandırır (dar marj,
aşırı öğrenme riski); küçük $C$ daha geniş marja izin verir (daha çok
düzenlileştirme).

\gsub{Kernel Hilesi (Kernel Trick)}

Veri doğrusal ayrılabilir değilse, SVM'i \emph{daha yüksek boyutlu} bir uzaya
taşırız; orada doğrusal ayrılabilir hale gelebilir. Kernel hilesi, bu dönüşümü
\emph{açıkça hesaplamadan}, yalnızca iç çarpımlar üzerinden yapar:
\[
K(\vek{x}, \vek{x}') = \phi(\vek{x})^\top\phi(\vek{x}').
\]

\begin{center}
\begin{tabularx}{\textwidth}{@{}>{\bfseries\color{anaMavi}}l >{$}l<{$} X@{}}
\toprule
\rowcolor{acikMavi}\color{anaMavi}\textbf{Kernel} & \color{anaMavi}\textbf{K(x,x')} & \color{anaMavi}\textbf{Kullanım}\\
\midrule
Lineer & \vek{x}^\top\vek{x}' & Doğrusal ayrılabilir, çok özellik\\
\rowcolor{acikGri} Polinom & (\gamma\vek{x}^\top\vek{x}'+r)^d & Polinom sınırlar\\
RBF (Gauss) & \exp(-\gamma\|\vek{x}-\vek{x}'\|^2) & En yaygın; esnek sınırlar\\
\bottomrule
\end{tabularx}
\end{center}

\begin{lstlisting}
from sklearn.svm import SVC
from sklearn.preprocessing import StandardScaler
from sklearn.pipeline import make_pipeline

# SVM ölçeğe duyarlıdır -> önce her zaman standartlaştır
model = make_pipeline(
    StandardScaler(),
    SVC(kernel='rbf', C=1.0, gamma='scale')   # RBF çekirdeği
)
model.fit(X_train, y_train)
print(model.score(X_test, y_test))
\end{lstlisting}

% =====================================================================
\gsec{Karar Ağaçları}
% =====================================================================

Karar ağacı, veriyi bir dizi ``evet/hayır'' sorusuyla bölerek sınıflandıran,
sezgisel ve yorumlanabilir bir modeldir. Her iç düğüm bir özellik üzerinde
bir eşik testi, her yaprak bir tahmindir. İnsanın karar verme sürecine
benzediği için ``beyaz kutu'' modeldir --- neden o kararı verdiğini
açıklayabilirsin.

\gsub{Bölme Kriteri: Bilgi Kazancı ve Gini}

Ağaç, her adımda veriyi \emph{en saf} alt gruplara bölen özelliği/eşiği seçer.
Saflık iki şekilde ölçülür. Kısım 1'deki entropiyi kullanan \emph{bilgi
kazancı}:
\[
IG = H(\text{ebeveyn}) - \sum_{k}\frac{m_k}{m}H(\text{çocuk}_k),
\]
veya hesaplaması daha ucuz olan \emph{Gini safsızlığı}:
\[
\text{Gini} = 1 - \sum_{c} p_c^2,
\]
burada $p_c$, düğümdeki $c$ sınıfının oranıdır. Her ikisi de düğüm ``saf''
(tek sınıf) olduğunda sıfırdır. Ağaç, bölme sonrası safsızlığı en çok azaltan
seçimi yapar.

\turetme{\textbf{Örnek:} 10 örnekli bir düğümde 5 pozitif, 5 negatif olsun.
Gini $= 1 - (0.5^2 + 0.5^2) = 0.5$ (maksimum karışıklık). Bir bölme, sol
düğüme 4 pozitif 1 negatif, sağa 1 pozitif 4 negatif koyarsa: her çocuğun
Gini'si $1-(0.8^2+0.2^2)=0.32$; ağırlıklı ortalama $0.32$. Safsızlık
$0.5\to0.32$ düştü --- iyi bir bölme.}

\begin{lstlisting}
from sklearn.tree import DecisionTreeClassifier, plot_tree

tree = DecisionTreeClassifier(
    criterion='gini',       # veya 'entropy'
    max_depth=4,            # aşırı öğrenmeyi önlemek için derinliği sınırla
    min_samples_leaf=5)
tree.fit(X_train, y_train)
print(tree.feature_importances_)    # en çok hangi özellikler önemli
\end{lstlisting}

\uyari{Karar ağaçları \emph{aşırı öğrenmeye çok eğilimlidir}: sınırlamazsan
her örneği ezberleyen dev bir ağaç kurar (eğitimde \%100, testte kötü). Bunu
\code{max\_depth}, \code{min\_samples\_leaf} gibi parametrelerle
``budarsın'' (pruning). Ama tek ağacın asıl gücü, aşağıdaki topluluk
yöntemlerinde ortaya çıkar.}

% =====================================================================
\gsec{Topluluk Yöntemleri (Ensemble)}
% =====================================================================

Topluluk yöntemleri, birçok ``zayıf'' modeli birleştirerek tek başına her
birinden daha güçlü bir model oluşturur. ``Kalabalığın bilgeliği'' ilkesi:
bağımsız hataları olan çok sayıda modelin ortalaması, bireysel hataları
söndürür. İki ana strateji vardır: bagging ve boosting.

\gsub{Bagging ve Random Forest}

\emph{Bagging} (Bootstrap Aggregating), eğitim verisinden rastgele
(yerine koyarak) alt örneklemler çeker, her birinde ayrı bir ağaç eğitir ve
tahminleri oylar/ortalar. Bu, variance'ı düşürür (aşırı öğrenmeyi azaltır).

\emph{Random Forest}, bagging'e bir ek yapar: her bölmede yalnızca özelliklerin
rastgele bir alt kümesini dener. Bu, ağaçları birbirinden \emph{ilişkisizleştirir}
--- ağaçlar ne kadar farklı hata yaparsa, ortalamaları o kadar iyi olur.

\begin{lstlisting}
from sklearn.ensemble import RandomForestClassifier

rf = RandomForestClassifier(
    n_estimators=200,       # ağaç sayısı
    max_depth=None,
    max_features='sqrt',    # her bölme için rastgele özellik alt kümesi
    random_state=42)
rf.fit(X_train, y_train)
print(rf.feature_importances_)
\end{lstlisting}

\gsub{Boosting: Hataları Ardışık Düzeltmek}

\emph{Boosting}, bagging'in aksine ağaçları \emph{paralel} değil \emph{ardışık}
kurar: her yeni model, önceki modellerin \emph{hatalarına} odaklanır. Gradient
Boosting bunu, her adımda önceki modelin \emph{artıklarına} (residual) yeni bir
ağaç uydurarak yapar --- adeta kaybın gradyanına gradient descent uygular.

\begin{center}
\begin{tabularx}{\textwidth}{@{}>{\bfseries\color{anaMavi}}l X X@{}}
\toprule
\rowcolor{acikMavi}\color{anaMavi}\textbf{} & \color{anaMavi}\textbf{Bagging (RF)} & \color{anaMavi}\textbf{Boosting}\\
\midrule
Ağaçlar & Paralel, bağımsız & Ardışık, birbirine bağlı\\
\rowcolor{acikGri} Azalttığı & Variance & Bias\\
Aşırı öğrenme & Dayanıklı & Daha hassas (dikkatli ayar)\\
\rowcolor{acikGri} Hız & Paralelleşir & Sıralı (yavaş)\\
\bottomrule
\end{tabularx}
\end{center}

\begin{lstlisting}
from sklearn.ensemble import GradientBoostingClassifier
# from xgboost import XGBClassifier          # popüler, daha hızlı alternatif

gb = GradientBoostingClassifier(
    n_estimators=100,
    learning_rate=0.1,      # her ağacın katkısını küçültür
    max_depth=3)            # boosting sığ "zayıf" ağaçlar kullanır
gb.fit(X_train, y_train)
\end{lstlisting}

\ipucu{Tablo (tabular) verisinde \textbf{gradient boosting} (XGBoost, LightGBM,
CatBoost) genelde en iyi performansı verir ve Kaggle yarışmalarının çoğunu
kazanır. Random Forest ise daha az ayar gerektiren, sağlam bir ilk tercihtir.
Derin öğrenme, tablo verisinde çoğu zaman bu ağaç tabanlı yöntemleri
geçemez --- asıl gücünü görüntü, ses ve metinde gösterir.}

% #####################################################################
\gpart{5}{Denetimsiz Öğrenme}
% #####################################################################

% =====================================================================
\gsec{K-Means Kümeleme}
% =====================================================================

Denetimsiz öğrenmede etiket yoktur; amaç veride gizli yapı bulmaktır. K-Means,
veriyi $k$ kümeye ayırır --- her kümeyi bir \emph{merkez} (centroid) temsil
eder ve her nokta en yakın merkeze atanır.

\gsub{Amaç Fonksiyonu}

K-Means, küme içi kareli uzaklıkların toplamını (inertia / WCSS) minimize eder:
\[
J = \sum_{k=1}^{K}\sum_{\vek{x}_i \in C_k} \|\vek{x}_i - \vek{\mu}_k\|_2^2,
\]
burada $\vek{\mu}_k$, $C_k$ kümesinin merkezidir. Bu problemi kesin çözmek
NP-zordur; bu yüzden \emph{Lloyd algoritması} ile yinelemeli olarak çözeriz.

\gsub{Lloyd Algoritması}

İki adım, yakınsayana dek tekrarlanır --- bu aslında bir tür ``koordinat inişi''
optimizasyonudur:

\begin{center}
\begin{tabularx}{\textwidth}{@{}>{\bfseries\color{anaMavi}}l X@{}}
\toprule
\rowcolor{acikMavi}\color{anaMavi}\textbf{Adım} & \color{anaMavi}\textbf{Yapılan}\\
\midrule
1. Atama & Her noktayı en yakın merkeze ata (merkezler sabit).\\
\rowcolor{acikGri} 2. Güncelleme & Her merkezi, atanan noktaların ortalaması yap (atamalar sabit).\\
\bottomrule
\end{tabularx}
\end{center}

\begin{lstlisting}
import numpy as np

class KMeansScratch:
    def __init__(self, k=3, max_iters=100, seed=42):
        self.k, self.max_iters, self.seed = k, max_iters, seed

    def fit(self, X):
        rng = np.random.default_rng(self.seed)
        # merkezleri k rastgele veri noktası olarak başlat
        idx = rng.choice(len(X), self.k, replace=False)
        self.centroids = X[idx].copy()

        for _ in range(self.max_iters):
            # Adım 1: her noktayı en yakın merkeze ata
            dists = np.linalg.norm(X[:, None] - self.centroids, axis=2)
            labels = np.argmin(dists, axis=1)

            # Adım 2: merkezleri atanan noktaların ortalaması olarak yeniden hesapla
            new_centroids = np.array([
                X[labels == j].mean(axis=0) if np.any(labels == j)
                else self.centroids[j]
                for j in range(self.k)
            ])
            if np.allclose(new_centroids, self.centroids):
                break                       # yakınsadı
            self.centroids = new_centroids
        self.labels_ = labels
        return self
\end{lstlisting}

\uyari{K-Means başlangıç merkezlerine \emph{duyarlıdır}: kötü başlangıç kötü
kümeye yakınsayabilir. \code{k-means++} başlatması bu riski büyük ölçüde
azaltır (scikit-learn varsayılanı budur). Ayrıca \emph{$k$'yı önceden bilmen
gerekir}; doğru $k$ için ``dirsek yöntemi'' (inertia'nın kırıldığı nokta) veya
``silhouette skoru'' kullanılır.}

\begin{lstlisting}
from sklearn.cluster import KMeans

km = KMeans(n_clusters=3, init='k-means++', n_init=10, random_state=42)
km.fit(X)
print(km.inertia_)          # küme içi kareler toplamı (J)
print(km.cluster_centers_)
\end{lstlisting}

% =====================================================================
\gsec{Temel Bileşen Analizi (PCA)}
% =====================================================================

PCA, çok boyutlu veriyi, \emph{bilginin çoğunu koruyarak} daha az boyuta
indirgeyen bir tekniktir. Fikir: verinin en çok \emph{yayıldığı} (en yüksek
varyanslı) yönleri bul; bu yönler verinin ``en önemli'' eksenleridir. Bu
yönlere \emph{temel bileşenler} denir.

\gsub{Matematiksel Türetme}

Veri $\vek{X}$'i önce \emph{merkezle} (her özellikten ortalamasını çıkar).
Merkezlenmiş veri için kovaryans matrisi:
\[
\vek{\Sigma} = \frac{1}{m}\vek{X}^\top\vek{X} \in \R^{n\times n}.
\]
Amacımız, üzerine izdüşürüldüğünde varyansı maksimize eden birim yön
$\vek{v}$'yi bulmaktır. İzdüşümün varyansı $\vek{v}^\top\vek{\Sigma}\vek{v}$'dir.

\turetme{$\|\vek{v}\|=1$ kısıtı altında $\vek{v}^\top\vek{\Sigma}\vek{v}$'yi
maksimize et. Lagrange çarpanı $\lambda$ ile:
\[
\mathcal{L}(\vek v,\lambda) = \vek{v}^\top\vek{\Sigma}\vek{v} - \lambda(\vek{v}^\top\vek{v}-1).
\]
$\vek v$'ye göre türev alıp sıfıra eşitleyince:
\[
2\vek{\Sigma}\vek{v} - 2\lambda\vek{v} = 0 \;\;\Longrightarrow\;\;
\boxed{\;\vek{\Sigma}\vek{v} = \lambda\vek{v}\;}.
\]
Bu tam olarak bir \emph{özdeğer denklemidir}! Yani en çok varyansa sahip yön,
kovaryans matrisinin \emph{en büyük özdeğerine} karşılık gelen özvektördür.
Üstelik o yöndeki varyans $\vek{v}^\top\vek\Sigma\vek v = \lambda$'ya eşittir
--- özdeğer, doğrudan o bileşenin açıkladığı varyanstır.}

Böylece PCA'nın reçetesi: kovaryans matrisinin özvektörlerini bul, özdeğerlere
göre büyükten küçüğe sırala, en büyük $d$ tanesini al --- bunlar aradığın
$d$ temel bileşendir.

\begin{lstlisting}
import numpy as np

class PCAScratch:
    def __init__(self, n_components):
        self.n_components = n_components

    def fit(self, X):
        self.mean_ = X.mean(axis=0)
        X_centered = X - self.mean_             # veriyi merkezle
        cov = np.cov(X_centered, rowvar=False)  # kovaryans matrisi (n, n)
        eigvals, eigvecs = np.linalg.eigh(cov)  # simetrik -> eigh kullan
        # özvektörleri özdeğere göre azalan sırada sırala
        order = np.argsort(eigvals)[::-1]
        self.components_ = eigvecs[:, order[:self.n_components]]
        self.explained_variance_ = eigvals[order[:self.n_components]]
        return self

    def transform(self, X):
        return (X - self.mean_) @ self.components_   # bileşenlere izdüşür
\end{lstlisting}

\begin{lstlisting}
from sklearn.decomposition import PCA
from sklearn.preprocessing import StandardScaler

X_scaled = StandardScaler().fit_transform(X)   # PCA ölçeklenmiş veri gerektirir
pca = PCA(n_components=2)
X_reduced = pca.fit_transform(X_scaled)
print(pca.explained_variance_ratio_)           # korunan varyans oranı
\end{lstlisting}

\ipucu{PCA \emph{ölçeğe duyarlıdır}: büyük değerli özellikler varyansa hâkim
olur. Bu yüzden PCA'dan önce mutlaka standartlaştır. \code{explained\_variance\_ratio\_}
toplamı, kaç bileşenle bilginin ne kadarını koruduğunu söyler --- genelde
\%95 varyansı koruyan bileşen sayısı seçilir.}

% #####################################################################
\gpart{6}{Sinir Ağları ve Derin Öğrenme}
% #####################################################################

% =====================================================================
\gsec{Perceptron'dan Çok Katmanlı Ağlara}
% =====================================================================

Sinir ağları, basit hesap birimlerini (nöronlar) katmanlar halinde
birleştirerek son derece karmaşık fonksiyonları öğrenebilen modellerdir.
Beynin nöronlarından ilham alır ama özünde saf matematiktir: matris çarpımları
ve doğrusal olmayan fonksiyonlar.

\gsub{Yapay Nöron}

Tek bir nöron, tıpkı lojistik regresyon gibi çalışır: girdilerin ağırlıklı
toplamını alır, bir \emph{aktivasyon fonksiyonundan} geçirir:
\[
a = g\Big(\sum_{j} w_j x_j + b\Big) = g(\vek{w}^\top\vek{x} + b).
\]
Tek nöron yalnızca \emph{doğrusal} sınırlar çizebilir. Gücün sırrı, nöronları
katmanlar halinde yığıp aralarına \emph{doğrusal olmayan} aktivasyonlar
koymaktır --- bu, ağın doğrusal olmayan fonksiyonları temsil etmesini sağlar.

\gsub{Aktivasyon Fonksiyonları}

Doğrusal olmama şarttır: aktivasyon olmadan, kaç katman olursa olsun ağ tek bir
doğrusal dönüşüme çöker. Yaygın seçenekler:

\begin{center}
\begin{tabularx}{\textwidth}{@{}>{\bfseries\color{anaMavi}}l >{$}l<{$} X@{}}
\toprule
\rowcolor{acikMavi}\color{anaMavi}\textbf{Fonksiyon} & \color{anaMavi}\textbf{g(z)} & \color{anaMavi}\textbf{Özellik}\\
\midrule
Sigmoid & \frac{1}{1+e^{-z}} & (0,1); doygunlukta gradyan kaybolur.\\
\rowcolor{acikGri} Tanh & \tanh(z) & (-1,1); sıfır merkezli.\\
ReLU & \max(0,z) & En yaygın; hızlı, gradyan kaybını azaltır.\\
\rowcolor{acikGri} Softmax & \frac{e^{z_i}}{\sum_j e^{z_j}} & Çok sınıflı çıkış (olasılıklar).\\
\bottomrule
\end{tabularx}
\end{center}

\bilgi{\textbf{Evrensel yaklaşım teoremi:} Yeterli sayıda nörona sahip tek
gizli katmanlı bir ağ, herhangi bir sürekli fonksiyonu istenen doğrulukta
yaklaşık temsil edebilir. Ama ``yapabilir'' ile ``öğrenebilir'' farklıdır ---
pratikte derin (çok katmanlı) ağlar, aynı işi çok daha az nöronla ve daha
kolay öğrenerek yapar. Derinliğin gücü buradadır.}

\gsub{İleri Besleme (Forward Propagation)}

Çok katmanlı bir ağda bilgi, girişten çıkışa katman katman akar. $l$. katman
için (matris biçiminde, tüm örnekler birden):
\[
\vek{Z}^{[l]} = \vek{A}^{[l-1]}\vek{W}^{[l]} + \vek{b}^{[l]}, \qquad
\vek{A}^{[l]} = g^{[l]}(\vek{Z}^{[l]}),
\]
burada $\vek{A}^{[0]}=\vek{X}$ (giriş), $\vek{A}^{[L]}=\hat{\vek{Y}}$ (çıkış).
Her katman bir doğrusal dönüşüm + doğrusal olmayan aktivasyondan ibarettir.

\begin{lstlisting}
import numpy as np

def relu(Z):
    return np.maximum(0, Z)

def softmax(Z):
    Z = Z - Z.max(axis=1, keepdims=True)        # sayısal kararlılık
    exp = np.exp(Z)
    return exp / exp.sum(axis=1, keepdims=True)

def forward(X, params):
    """İki katmanlı ağ: X -> ReLU -> softmax."""
    Z1 = X @ params['W1'] + params['b1']
    A1 = relu(Z1)
    Z2 = A1 @ params['W2'] + params['b2']
    A2 = softmax(Z2)
    cache = (X, Z1, A1, Z2, A2)                  # backprop için saklandı
    return A2, cache
\end{lstlisting}

% =====================================================================
\gsec{Geri Yayılım (Backpropagation)}
% =====================================================================

Geri yayılım, sinir ağlarının nasıl öğrendiğinin kalbidir: kaybın, ağdaki
\emph{her ağırlığa} göre gradyanını verimli biçimde hesaplar. Özü, Kısım 1'deki
\emph{zincir kuralıdır} --- hata, çıkıştan girişe doğru katman katman ``geri
yayılır''. Bu bölüm, iki katmanlı bir ağ için gradyanları baştan sona türetir.

\gsub{Neden Zincir Kuralı?}

Kayıp $J$, çıkışa; çıkış, son katman ağırlıklarına; o ağırlıklar bir önceki
katmanın çıkışına\dots\ bağlıdır. Bir ara ağırlığın kaybı nasıl etkilediğini
bulmak için zincir kuralıyla bu bağımlılık zincirini çarparız:
\[
\frac{\partial J}{\partial \vek{W}^{[l]}} =
\frac{\partial J}{\partial \vek{A}^{[l]}}\cdot
\frac{\partial \vek{A}^{[l]}}{\partial \vek{Z}^{[l]}}\cdot
\frac{\partial \vek{Z}^{[l]}}{\partial \vek{W}^{[l]}}.
\]

\gsub{Tam Türetme (Softmax + Çapraz Entropi)}

\turetme{Ağ: $\vek{X}\to\vek{Z}^{[1]}\to\vek{A}^{[1]}=\text{ReLU}\to
\vek{Z}^{[2]}\to\vek{A}^{[2]}=\text{softmax}$, kayıp çapraz entropi.
Türetmeyi çıkıştan geriye yaparız.

\textbf{(1) Çıkış katmanı.} Softmax + çapraz entropinin birlikte türevi
olağanüstü sadeleşir (bu ikilinin birlikte kullanılmasının nedeni budur):
\[
\frac{\partial J}{\partial \vek{Z}^{[2]}} = \vek{A}^{[2]} - \vek{Y}
\;\equiv\; \vek{\delta}^{[2]}.
\]
Buradan son katman gradyanları:
\[
\frac{\partial J}{\partial \vek{W}^{[2]}} = \frac{1}{m}(\vek{A}^{[1]})^\top\vek{\delta}^{[2]},
\qquad
\frac{\partial J}{\partial \vek{b}^{[2]}} = \frac{1}{m}\sum_i \vek{\delta}^{[2]}_i.
\]

\textbf{(2) Hatayı geri yay.} Gizli katmanın hatası, çıkış hatasının ağırlıklarla
geri taşınıp ReLU türeviyle çarpılmasıdır:
\[
\vek{\delta}^{[1]} = \big(\vek{\delta}^{[2]}(\vek{W}^{[2]})^\top\big)\odot g'(\vek{Z}^{[1]}),
\qquad g'(z)=\mathbb{1}[z>0]\ \ (\text{ReLU türevi}).
\]
Burada $\odot$ eleman bazlı çarpımdır.

\textbf{(3) İlk katman gradyanları.} Aynı kalıp:
\[
\frac{\partial J}{\partial \vek{W}^{[1]}} = \frac{1}{m}\vek{X}^\top\vek{\delta}^{[1]},
\qquad
\frac{\partial J}{\partial \vek{b}^{[1]}} = \frac{1}{m}\sum_i \vek{\delta}^{[1]}_i.
\]}

\begin{lstlisting}
def backward(Y, cache, params):
    """Yukarıdaki iki katmanlı ağ için geri yayılım.
    Y: one-hot kodlanmış gerçek etiketler (m, n_classes)."""
    X, Z1, A1, Z2, A2 = cache
    m = X.shape[0]
    grads = {}

    # (1) Çıkış katmanı hatası: softmax + çapraz entropi => A2 - Y
    dZ2 = A2 - Y                                  # (m, n_classes)
    grads['W2'] = (A1.T @ dZ2) / m
    grads['b2'] = dZ2.sum(axis=0, keepdims=True) / m

    # (2) Hatayı gizli katmana geri yay (ReLU türevi)
    dA1 = dZ2 @ params['W2'].T
    dZ1 = dA1 * (Z1 > 0)                           # ReLU'(z) = 1[z>0]

    # (3) İlk katman gradyanları
    grads['W1'] = (X.T @ dZ1) / m
    grads['b1'] = dZ1.sum(axis=0, keepdims=True) / m
    return grads

def update(params, grads, lr):
    for key in params:
        params[key] -= lr * grads[key]            # gradyan inişi adımı
    return params
\end{lstlisting}

\bilgi{Backprop'un ``geri yayma'' kalıbına dikkat et: her katmanın gradyanı
$\vek{A}^\top\vek{\delta}$ biçimindedir ve hata $\vek{\delta}$, bir üst
katmandan $\vek{W}^\top$ ile geri taşınır. Bu tekrarlayan yapı, PyTorch/TensorFlow
gibi kütüphanelerin \emph{otomatik türev} mekanizmasının temelidir --- sen
ileri geçişi yaz, onlar bu zinciri otomatik kurar.}

\gsub{Kaybolan/Patlayan Gradyan}

Derin ağlarda gradyan, katmanlar boyunca çarpıla çarpıla geri yayılır. Çarpanlar
sürekli $<1$ ise gradyan katman katman \emph{küçülüp kaybolur} (vanishing);
$>1$ ise \emph{patlar} (exploding). Sigmoid/tanh doygunlukta gradyanı
sıfıra iter --- bu yüzden derin ağlarda ReLU tercih edilir. Ayrıca dikkatli
ağırlık başlatma (Xavier/He), batch normalization ve residual bağlantılar
bu sorunu hafifletir.

% =====================================================================
\gsec{Optimizasyon Algoritmaları}
% =====================================================================

Düz gradient descent, sinir ağlarını eğitmek için genelde yavaştır ve dar
vadilerde salınır. Modern optimize ediciler, geçmiş gradyanları kullanarak
öğrenmeyi hızlandırır ve kararlı hale getirir.

\gsub{Momentum}

Momentum, gradyanların bir \emph{hareketli ortalamasını} (hız) tutar; böylece
tutarlı yönlerde ivmelenir, gürültülü salınımları söndürür --- yokuş aşağı yuvarlanan
bir topa benzer:
\[
\vek{v} \leftarrow \beta\vek{v} + (1-\beta)\nabla J, \qquad
\vek{w} \leftarrow \vek{w} - \eta\vek{v}.
\]

\gsub{Adam: En Yaygın Optimize Edici}

Adam (Adaptive Moment Estimation), momentum ile parametre-başına uyarlamalı
öğrenme oranını birleştirir. Gradyanın hem birinci momentini (ortalama, $\vek m$)
hem ikinci momentini (varyans, $\vek v$) tutar:
\[
\vek{m} \leftarrow \beta_1\vek{m} + (1-\beta_1)\vek{g},\qquad
\vek{v} \leftarrow \beta_2\vek{v} + (1-\beta_2)\vek{g}^2,
\]
\[
\hat{\vek m}=\frac{\vek m}{1-\beta_1^t},\quad
\hat{\vek v}=\frac{\vek v}{1-\beta_2^t},\qquad
\vek{w}\leftarrow\vek{w}-\eta\frac{\hat{\vek m}}{\sqrt{\hat{\vek v}}+\epsilon}.
\]
$\hat{\vek m},\hat{\vek v}$ ``bias düzeltmesidir'' (başlangıçta sıfıra yakın
olmayı telafi eder). Adam çoğu problemde iyi çalışır ve az ayar gerektirir ---
bu yüzden derin öğrenmede varsayılan tercihtir.

\begin{center}
\begin{tabularx}{\textwidth}{@{}>{\bfseries\color{anaMavi}}l X X@{}}
\toprule
\rowcolor{acikMavi}\color{anaMavi}\textbf{Optimize edici} & \color{anaMavi}\textbf{Fikir} & \color{anaMavi}\textbf{Not}\\
\midrule
SGD & Saf gradyan adımı & Basit, ayar ister\\
\rowcolor{acikGri} Momentum & Geçmiş gradyanla ivmelen & Salınımı azaltır\\
RMSProp & Uyarlamalı öğrenme oranı & RNN'lerde iyi\\
\rowcolor{acikGri} Adam & Momentum + RMSProp & En yaygın varsayılan\\
\bottomrule
\end{tabularx}
\end{center}

% =====================================================================
\gsec{Derin Öğrenme Mimarileri ve Pratik}
% =====================================================================

Tam bağlı ağlar (MLP) her şeyi öğrenebilir ama görüntü, dizi gibi \emph{yapılı}
veride verimsizdir. Özel mimariler, verinin yapısını mimariye gömer.

\gsub{Evrişimli Ağlar (CNN)}

CNN'ler görüntü için tasarlanmıştır. Tam bağlantı yerine, küçük \emph{filtreleri}
(kernel) görüntü üzerinde kaydırarak yerel örüntüleri (kenar, köşe, doku)
yakalar. Üç temel fikir: \textbf{yerel bağlantı} (her nöron sadece küçük bir
bölgeye bakar), \textbf{ağırlık paylaşımı} (aynı filtre her yerde kullanılır ---
çok az parametre) ve \textbf{havuzlama} (pooling ile boyut indirgeme). Katmanlar
derinleştikçe basit örüntülerden karmaşık nesnelere doğru hiyerarşik özellikler
öğrenilir.

\gsub{Tekrarlayan Ağlar (RNN) ve Transformer}

RNN'ler dizisel veri (metin, zaman serisi, ses) içindir: bir \emph{gizli durum}
tutarak geçmiş bilgiyi taşırlar. Uzun bağımlılıklarda kaybolan gradyan sorunu
yaşarlar; \textbf{LSTM} ve \textbf{GRU} kapı (gate) mekanizmalarıyla bunu çözer.
Günümüzde metin ve giderek daha çok alanda, dizinin tüm parçalarına aynı anda
bakan \emph{dikkat} (attention) mekanizmasına dayalı \textbf{Transformer}
mimarisi baskındır --- büyük dil modellerinin (LLM) temelidir.

\begin{lstlisting}
# Modern derin öğrenme sıfırdan değil, framework'lerle yapılır.
# Yukarıda türettiğimiz iki katmanlı ağın minimal bir PyTorch örneği:
import torch
import torch.nn as nn

model = nn.Sequential(
    nn.Linear(784, 128),      # giriş -> gizli
    nn.ReLU(),
    nn.Linear(128, 10)        # gizli -> çıkış (10 sınıf)
)
criterion = nn.CrossEntropyLoss()          # softmax + çapraz entropi
optimizer = torch.optim.Adam(model.parameters(), lr=1e-3)

for epoch in range(n_epochs):
    for xb, yb in dataloader:              # mini-gruplar
        optimizer.zero_grad()              # gradyanları sıfırla
        preds = model(xb)                  # ileri geçiş
        loss = criterion(preds, yb)
        loss.backward()                    # backprop (otomatik!)
        optimizer.step()                   # Adam güncellemesi
\end{lstlisting}

\ipucu{Otomatik türev sayesinde artık gradyanları elle yazmıyoruz ---
\code{loss.backward()} tüm zinciri otomatik hesaplar. Ama bu rehberde backprop'u
elle türetmemizin nedeni tam da bu: ``sihrin'' altında sadece zincir kuralı
olduğunu bilmek, ağları hata ayıklamak, yeni katmanlar tasarlamak ve neyin
neden çalışmadığını anlamak için vazgeçilmezdir.}

% =====================================================================
\gsec{Kapanış: Bir Bütün Olarak Makine Öğrenmesi}
% =====================================================================

Bu rehber boyunca gördük ki makine öğrenmesi, üç matematiksel temelin
kaynaşmasıdır:

\begin{center}
\begin{tabularx}{\textwidth}{@{}>{\bfseries\color{anaMavi}}l X@{}}
\toprule
\rowcolor{acikMavi}\color{anaMavi}\textbf{Temel} & \color{anaMavi}\textbf{Rolü}\\
\midrule
Lineer cebir & Veriyi ve modeli temsil eder (vektör/matris işlemleri).\\
\rowcolor{acikGri} Kalkülüs \& optimizasyon & Kaybı minimize ederek modeli eğitir (gradient descent).\\
Olasılık \& istatistik & Belirsizliği modeller, kayıpları türetir (MLE, cross-entropy).\\
\bottomrule
\end{tabularx}
\end{center}

Aynı fikirlerin tekrar tekrar döndüğüne dikkat et: gradyan hem lineer
regresyonda, hem lojistik regresyonda, hem sinir ağında \emph{aynı biçimi}
alır ($\vek{X}^\top(\hat{\vek y}-\vek y)$). MLE hem MSE'yi hem çapraz entropiyi
doğurur. Bir kez bu ortak yapıyı gördüğünde, yeni bir algoritma artık ezberlenecek
bir formül değil, aynı temaların bir varyasyonu olur.

\gsub{Bir Model Nasıl Seçilir?}

\begin{center}
\begin{tabularx}{\textwidth}{@{}>{\bfseries\color{anaMavi}}l X@{}}
\toprule
\rowcolor{acikMavi}\color{anaMavi}\textbf{Durum} & \color{anaMavi}\textbf{Önerilen model}\\
\midrule
Az veri, yorumlanabilirlik & Lineer/lojistik regresyon, karar ağacı\\
\rowcolor{acikGri} Tablo veri, en iyi performans & Gradient boosting (XGBoost/LightGBM)\\
Sağlam, az ayarlı ilk deneme & Random Forest\\
\rowcolor{acikGri} Görüntü & CNN (evrişimli ağ)\\
Metin / dizi & Transformer (veya LSTM)\\
\rowcolor{acikGri} Etiketsiz, grup bulma & K-Means, PCA\\
\bottomrule
\end{tabularx}
\end{center}

\ipucu{Son bir tavsiye: \textbf{basit başla}. Her zaman basit bir temel çizgi
(baseline) modelle başla --- lineer/lojistik regresyon veya random forest.
Karmaşık bir modele ancak basit olan yetmediğinde geç. Çoğu zaman iyi
\emph{özellik mühendisliği} ve temiz veri, fancy bir modelden çok daha fazla
kazandırır. Ve her şeyi bir doğrulama kümesinde ölç --- eğitim skoruna asla
güvenme. İyi öğrenmeler!}

\vfill
\begin{center}
{\color{ortaGri}\rule{0.5\linewidth}{0.6pt}}\\[6pt]
{\color{ortaGri}\small Rehberin sonu --- Matematikten koda, iyi öğrenmeler!}
\end{center}

\end{document}
